% 本章包含长脚注分页控制命令；主文件导言区需加载 needspace，
% 并设置 \interfootnotelinepenalty=10000（见随文说明）。
\part*{微模块1\quad 人工智能通识应用}
\addcontentsline{toc}{part}{微模块1\texorpdfstring{\quad}{ }人工智能通识应用}

% ============================================================
% 第4章：人工智能通识应用
% ============================================================

\setcounter{chapter}{3}

\chapter{人工智能通识应用}

% 表题注采用“表4-1 题名”的分章编号格式
\renewcommand{\thetable}{\thechapter-\arabic{table}}
\captionsetup[table]{labelformat=empty,labelsep=none}

% 网络来源的著录信息与访问日期放在正文对应页的脚注中
\newcommand{\chaptertablelabel}{{\xeCJKsetup{CJKecglue={}}表\thetable}}

\section*{本章导览}

本章围绕人工智能在大学生日常学习、创作表达和公共发布中的通识应用展开。这里的通识应用，不是要求学生掌握复杂算法，也不是把 AI 当作万能工具，而是学习在真实任务中合理使用 AI：先明确任务目标，再提供必要素材，再获得初稿或方案，生成后进行人工核对、修改和发布前审查。通过这些训练，学生能够理解 AI 的价值不只在于生成内容，更在于帮助人整理信息、扩展思路、优化表达和提高工作效率。

本章内容按照学生常见使用场景组织，包括正式材料写作、办公材料整理、创意娱乐生成和平台内容发布四类任务。正式材料写作关注简历、申请信和感谢信；办公材料整理与文档写作关注材料提取、活动总结、方案撰写和公文纪要；创意娱乐生成关注小说、漫画和照片修复；平台发布关注社交笔记、公众号推文和多平台改写。四类场景共同指向一个目标：让学生学会在真实任务中合理使用 AI，并在生成后进行必要审查。

应用场景虽然形式不同，但使用 AI 的基本方法是相通的。使用 AI 时，学生需要把任务说清楚，包括使用对象、写作目的、素材范围、输出格式和语气要求；也要把事实边界说清楚，明确哪些信息可以整理、哪些信息不能编造、哪些内容需要保留原貌；还要把审查责任留给使用者，不能因为 AI 输出流畅就直接采用。尤其在正式材料、公文写作、照片修复和平台发布中，真实性、准确性、隐私保护和授权意识都比语言是否好看更重要。

学习本章时，学生应重点关注每个案例中的输入方式、输出差异和修改过程。简单提问往往只能得到通用模板，提供真实素材和明确要求后，AI 才能生成更可用的内容；生成结果看似完整，也仍然需要人工检查事实、格式、语气、逻辑和使用边界。通过本章学习，学生应逐步形成一种基本能力：能够把 AI 当作辅助工具，而不是替代判断的工具；能够使用 AI 提高表达效率，也能够识别并修正 AI 输出中的不准确、不真实和不合适之处。

\newpage

\section{简历材料生成与信件写作}

\begin{sectionkeypoints}
\item 理解正式材料写作的基本任务，认识简历、申请信、自我陈述和感谢信等文本在展示经历、说明理由和表达态度中的使用要求。
\item 掌握使用 AI 辅助生成正式材料的基本思路，能够围绕用途、对象、真实经历、格式和语气要求组织输入内容。
\item 了解简历、申请信和感谢信的表达重点，理解经历筛选、申请目标说明、具体感谢内容和后续态度之间的区别。
\item 熟悉 AI 初稿修改与审查的基本方法，能够检查事实真实性、表述具体性、语气得体性和格式清晰度。
\item 认识正式材料写作的使用边界，明确不能虚构经历、夸大能力、替他人作出不确定评价或暴露无关个人信息。
\end{sectionkeypoints}

正式材料写作是大学生在学习发展、组织参与和日常交流中经常遇到的基础任务。它不同于一般聊天，也不同于随意记录，而是面向明确对象、服务具体目的、需要承担一定责任的书面表达。学生在申请社团、竞选学生组织、报名项目、参加实习、联系老师和表达感谢时，往往都需要准备相对正式的材料。这类材料看起来篇幅不一定很长，但它直接影响他人对个人经历、表达能力和做事态度的判断。

本节围绕正式材料写作展开，重点不是讲文体概念，而是讲如何把 AI 用到真实任务中。正式材料创作通常包含三个连续环节：先把个人经历和真实信息整理出来，再根据使用场景生成合适的文本初稿，最后对 AI 输出内容进行打磨和审查，确保材料真实、具体、得体。AI 可以帮助学生提高整理材料和形成初稿的效率，但不能替学生编造经历，也不能替学生承担最终判断责任。

为了让学习过程更连贯，本节讲述了一个完整案例：小李准备竞选学生会新媒体部干事。竞选前，她需要先制作简历，让老师和学生会负责人快速了解自己的专业背景、相关经历和能力基础；随后，她需要提交申请信，说明自己为什么想加入新媒体部、过去有哪些准备、入选后希望承担哪些工作；竞选成功后，她还需要写感谢信，向面试的学长学姐表达感谢，并说明自己后续会认真参与部门工作。简历、申请信和感谢信分别对应正式材料写作中的展示经历、表达意愿和维护关系三个任务。

从这个案例可以看到，使用 AI 写正式材料并不是简单输入一句帮我写，然后直接复制结果。更合理的做法是先明确用途，再整理材料，再让 AI 生成初稿，生成后由使用者进行人工修改。简历要突出真实经历和岗位相关性，申请信要说明动机和匹配度，感谢信要体现具体感谢原因和后续态度。三类材料虽然形式不同，但共同要求都是：事实不能虚构，表达不能夸张，语气要符合场景，个人信息要适度提供。

因此，本节按照任务发生的顺序展开。本节先学习如何使用 AI 制作简历，解决经历如何筛选和呈现的问题；再学习如何使用 AI 写申请信，解决申请理由如何表达清楚的问题；还会学习如何使用 AI 写感谢信，解决感谢内容如何自然得体的问题。通过这三个任务，学生可以掌握 AI 辅助正式材料写作的一般方法，也能理解正式材料创作中必须坚持的真实性、准确性和场景意识。


\subsection{简历材料整理与经历呈现}

简历用于快速呈现个人经历和能力基础。本小节以小李竞选学生会新媒体部干事为案例，先通过简单提问观察 AI 输出的局限，再补充简历创作的前置知识和真实材料，再完成简历成稿与审查。

在大学生活中，简历并不只出现在求职场景中。学生参加社团竞选、报名学生组织、申请实习岗位、寻找竞赛队友、参加项目申报，甚至在某些课程展示中进行个人介绍时，都可能需要一份简洁清楚的个人材料。简历的作用，是让阅读者在较短时间内了解一个人的基本情况、经历基础和能力特点。它不是个人经历的简单堆积，也不是把所有优点都写上去，而是根据具体用途，选择最相关的内容进行呈现。

AI 可以在简历写作中发挥辅助作用。它可以帮助学生梳理经历、归类信息、修改表达、压缩语言，也可以帮助学生把零散材料整理成较清晰的简历结构。但学生需要明白，AI 不是用来编造一份好看的简历的工具。真正有效的简历，必须建立在真实经历的基础上。AI 能做的是帮助整理和表达，而不是替学生制造经历。

\subsubsection*{\textbf{• 案例}}

小李是数字媒体技术专业大二学生，准备竞选学生会新媒体部干事。在正式提交申请前，她需要先准备一份简历，用较短的篇幅呈现自己的专业背景、相关经历、技能基础和竞选方向。对于老师和学生会负责人来说，简历的作用不是展开叙述所有经历，而是快速判断小李是否具备新媒体部工作所需要的图文排版、素材整理、短视频制作和活动宣传意识。因此，在这一阶段，小李可以使用 AI 辅助梳理已有经历，提炼与竞选岗位相关的信息，并生成一份结构清晰、重点突出的简历。


小李最初并不清楚要向 AI 提供哪些材料。她只是想先看看 AI 会生成什么，于是进行了一个很简单的提问。

\begin{userbox}
帮小李写一份竞选学生会新媒体部干事的简历。
\end{userbox}

\begin{aibox}
\textbf{姓名：}小李

\textbf{竞选岗位：}学生会新媒体部干事

\textbf{个人简介：}本人希望竞选学生会新媒体部干事，对校园宣传和新媒体工作有兴趣，愿意在学生会工作中学习和锻炼。

\textbf{相关经历：}暂未提供具体经历。建议补充与推文编辑、图片处理、视频制作、活动宣传、团队协作等相关的真实经历。

\textbf{技能基础：}暂未提供具体技能。建议补充会使用的软件、完成过的作品或参与过的活动。

\textbf{自我评价：}本人愿意学习新媒体相关工作，能够认真对待竞选和后续任务。具体评价可结合真实经历进一步完善。
\end{aibox}

这次回答和提问之间是对应的。由于小李只告诉 AI 竞选学生会新媒体部干事，AI 只能确定姓名和竞选岗位，不能准确写出她会什么、做过什么。如果 AI 在没有材料的情况下直接写出熟练掌握剪映\footnote{剪映是一款视频剪辑软件，仅做示例。}，参与过公众号排版、负责过活动摄影等内容，这些内容就可能是编造的。教材中设置这个对比例子，是为了提醒学生：简短提问不是不能用，但简短提问通常只能得到框架或模板，不能得到真正可信的个人材料。

\subsubsection*{\textbf{• 内容整理}}

使用 AI 生成简历时，学生要先把简历任务说明清楚，包括简历用途、投递对象、申请岗位或竞选方向，以及希望重点呈现的能力。简历生成不能只依靠一句“帮我写简历”，而应让 AI 明确这份材料服务于什么场景，从而围绕目标筛选内容。

同时，要向 AI 提供真实、具体的个人材料。可以包括基本信息、专业年级、相关经历、技能基础、实践成果和个人意愿等内容。提供经历时，应尽量写清楚做过什么、参与到什么程度、和目标岗位有什么关系，避免只提供性格开朗、积极上进等空泛评价。

还要说明输出结构和表达限制。学生可以要求 AI 按基本信息、申请方向、相关经历、技能基础、自我评价等部分整理，语言简洁正式，不要夸大能力，不要补充未提供的获奖、职务或项目经历。整理好这些内容后，小李把自己的真实材料输入给 AI，并提出了更加明确的生成要求。

\begin{userbox}
请帮我为小李整理一份用于竞选学生会新媒体部干事的简历。不要编造内容，只根据我提供的信息写。

小李是数字媒体技术专业大二学生，想竞选学生会新媒体部干事。她曾参与班级公众号图文排版，协助整理过三篇班级活动推文初稿；在校运动会期间参与活动照片拍摄和分类整理；平时会使用剪映制作短视频，也会用 Canva\footnotemark 制作简单海报。她希望加入新媒体部后，参与校园活动宣传、推文编辑、素材整理和短视频制作。

请按照基本信息、竞选方向、相关经历、技能基础、自我评价的结构生成。语言要简洁正式，不要写得过度夸张。
\end{userbox}

\footnotetext{Canva 是一款在线视觉设计工具，仅做示例。}

\begin{aibox}
\textbf{基本信息}

姓名：小李

专业年级：数字媒体技术专业大二学生

\textbf{竞选方向}

学生会新媒体部干事

\textbf{相关经历}

1. \textbf{班级公众号图文排版经历}：曾参与班级公众号图文排版，协助整理三篇班级活动推文初稿，了解推文内容整理和版面呈现的基本流程。

2. \textbf{校园活动素材整理经历}：在校运动会期间参与活动照片拍摄和分类整理，能够根据活动内容对图片素材进行初步筛选和归类。

3. \textbf{短视频与海报制作经历}：平时会使用剪映制作短视频，也会用 Canva 制作简单海报，具备一定的新媒体内容制作基础。

\textbf{技能基础}

具备基础图文排版、活动照片整理、短视频剪辑和简单海报制作能力；对校园活动宣传、推文编辑和短视频表达有较强兴趣。

\textbf{自我评价}

我对新媒体内容制作较感兴趣，也愿意在学生会工作中继续提升图文编辑、素材整理、视频剪辑和团队协作能力。如果能够加入新媒体部，我会认真参与校园活动宣传相关工作，按时完成部门安排的任务。
\end{aibox}

第二次生成的简历明显更具体。它没有凭空增加小李没有提供的经历，而是围绕已有材料进行整理。比如“协助整理三篇班级活动推文初稿”“参与活动照片拍摄和分类整理”“会使用剪映和 Canva”等内容，都能在提问中找到来源。这样的输出才适合作为简历初稿。

\subsubsection*{\textbf{• 成稿示例}}

经过 AI 初稿整理后，小李还需要做进一步修改。简历不是越长越好。对于学生会干事竞选来说，一页左右的简历更适合阅读。因此，小李可以继续使用 AI 把较长的句子压缩，把重复表达删去，形成更清晰的版本，并根据信息进行简历生成，生成的版式示例见图\ref{fig:ch4_01_resume_example}。

这个成稿版本比 AI 初稿更适合实际使用。它保留了真实经历，也减少了重复叙述。学生在写简历时，可以把 AI 作为整理助手，先让 AI 帮助分类，再由自己根据场景决定保留哪些内容。

\par
\noindent\begin{minipage}{\textwidth}
\centering
\BookFigureImage{./images_chapter_04/简历图片.png}
\captionof{figure}{AI 生成简历示例图}
\label{fig:ch4_01_resume_example}
\end{minipage}
\par

\par\vspace{1em}



前面的简历示例已经形成了可提交的初稿，但小李还不能只看版面是否整齐。正式使用前，需要把案例中的经历表述、能力描述和个人信息逐项核对。

\Needspace{12\baselineskip}
\begin{critiquebox}
简历审查应先抓三件事：经历真实、内容相关、信息必要。小李可以写“协助整理三篇推文初稿”，不能写成“独立负责公众号运营”；可以写“会用 Canva 制作简单海报”，不宜写成“精通平面设计”。正式简历以真实呈现为底线\footnotemark。

内容要服务投递或竞选目标。新媒体部竞选应优先保留排版、照片整理、短视频制作等经历，关联不强的活动可以压缩。个人信息也要控制在必要范围内，身份证号、家庭住址、完整学号等不宜输入 AI 工具，联系方式是否保留应看提交对象和使用场景\footnotemark。

\textbf{审查要点：}是否把参与写成负责，是否堆入无关经历，是否用空泛评价代替具体任务，是否暴露过多个人信息，是否把能力表述写得过满。
\end{critiquebox}
\addtocounter{footnote}{-2}
\stepcounter{footnote}\footnotetext{参见：王晨、刘则成《大学生如何进行简历制作》，中国大学生在线，2024年10月11日。该文提醒简历制作应以真诚为出发点，真实呈现个人情况。访问日期：2026年7月18日。}
\stepcounter{footnote}\footnotetext{参见：全国人民代表大会常务委员会《中华人民共和国个人信息保护法》（2021年8月20日公布，2021年11月1日起施行）第五条、第六条。访问日期：2026年7月18日。}



通过这一过程，小李完成了第一份正式材料。简历让评审者能够快速了解她有什么经历和基础，也为后续申请信提供了具体材料。


\subsection{申请理由组织与自我陈述}

申请信和自我陈述的核心作用，是说明申请者为什么争取某个机会，以及自己为什么适合这个机会。本小节围绕申请信创作展开：先说明申请信在正式材料写作中的作用和常见使用场景，再结合小李的竞选材料，展示如何把简历中的经历转化为连贯的申请理由，并说明申请信完成后需要从内容真实、理由充分和语气得体等方面进行审查。

申请信和自我陈述常用于争取机会的场景。学生申请奖学金、实习岗位、社团部门、学生组织、科研项目、竞赛团队时，往往需要用文字说明自己的申请目的、已有基础和未来打算。不同场景的写作重点有所不同：奖学金申请更强调学习表现和综合素质，实习申请更强调专业背景和岗位匹配度，社团申请更强调兴趣基础和参与意愿，项目申报更强调已有准备和预期贡献。虽然场景不同，但基本逻辑相似，都是先说明申请目的，再说明已有基础，最后说明未来打算。

与简历相比，申请信更强调表达的连贯性和说服力。简历回答的是我有什么经历，申请信回答的是我为什么申请和我为什么适合。因此，申请信不能只是把简历内容换一种说法重复一遍，而要把个人经历、申请目标和未来计划连接起来，让阅读者看到申请者的动机、基础和态度。

\subsubsection*{\textbf{• 案例}}

小李已经整理好了用于竞选学生会新媒体部干事的简历，里面包含她的专业背景、新媒体相关经历和技能基础。接下来，她需要把这些材料转化为一封申请信，让学生会负责人看到她的申请动机和工作意愿。AI 可以帮助学生搭建申请信结构，也可以将简历中的条目转化为连贯表达。但如果提问信息太少，AI 输出通常会比较空泛。小李一开始只进行了较为简单的提问。

\begin{userbox}
帮小李写一封竞选学生会新媒体部干事的申请信。
\end{userbox}

\begin{aibox}
\noindent 尊敬的学生会负责人：

\par\qquad 您好！

\par\qquad 我是小李，希望竞选学生会新媒体部干事。我对新媒体工作感兴趣，也希望通过学生会这个平台锻炼自己、服务同学。如果能够加入新媒体部，我会认真学习相关工作，积极配合部门安排，努力完成交给我的任务。

\par\qquad 恳请学生会给予我一次学习和实践的机会。感谢您的阅读！

\par\qquad 此致

\noindent 敬礼！

\par\hfill 申请人：小李

\par\hfill xx年xx月xx日
\end{aibox}

这段输出没有明显错误，但内容比较普通。它只使用了提问中已有的信息：小李、竞选学生会新媒体部干事。由于提问没有提供小李的专业、经历和技能，AI 无法写出更具体的内容。这样的申请信可以作为框架，但不能直接作为正式提交材料。

\subsubsection*{\textbf{• 内容整理}}

使用 AI 生成申请信时，学生应先说明申请任务，明确申请对象和申请目标，例如申请加入某个部门、参加某个项目或争取某个机会；接着补充申请原因，包括个人兴趣、申请动机和对目标岗位或活动的基本理解；还需要提供与申请目标相关的真实经历和能力基础，让 AI 根据已有材料展开，而不是生成空泛表态。表达要求也要一并说明，例如语气正式自然、内容真实具体、结构清楚，并提醒 AI 不要夸大经历、不要补充未提供的信息。这样生成的申请信才能更符合正式材料写作的要求。

小李在修改提问时，就按照这个思路整理材料。她没有让 AI 随意发挥，而是把申请目的、已有基础和加入后的打算都写清楚。

\begin{userbox}
请根据以下信息，为小李写一封竞选学生会新媒体部干事的申请信。要求语气正式、自然，不要写得过度夸张，字数 500 字左右。

小李是数字媒体技术专业大二学生，想竞选学生会新媒体部干事。她曾参与班级公众号图文排版，协助整理过三篇班级活动推文初稿；在校运动会期间参与活动照片拍摄和分类整理；平时会使用剪映制作短视频，也会用 Canva 制作简单海报。她希望加入新媒体部后，参与校园活动宣传、推文编辑、素材整理和短视频制作。

申请信需要写清楚：申请目的、已有基础、加入后的打算和结尾表达。
\end{userbox}

\begin{aibox}
\noindent 尊敬的学生会负责人：

\par\qquad 您好！

\par\qquad 我是数字媒体技术专业的大二学生小李，希望能够竞选学生会新媒体部干事。进入大学以来，我一直关注校园活动的宣传与记录，也希望通过学生会新媒体部这个平台，进一步参与图文编辑、素材整理和短视频制作等实践工作。

\par\qquad 在以往的学习和活动中，我接触过一些新媒体相关内容。我曾参与班级公众号图文排版，协助整理过三篇班级活动推文初稿，对推文内容整理和版面呈现有一定了解。在校运动会期间，我参与活动照片拍摄和分类整理，认识到校园宣传不仅需要记录现场，也需要根据主题筛选合适素材。平时我也会使用剪映制作短视频，并尝试用 Canva 制作简单海报。这些经历让我对新媒体内容制作有了更明确的兴趣。

\par\qquad 如果能够加入学生会新媒体部，我希望参与校园活动宣传、推文编辑、素材整理和短视频制作等工作。在完成部门任务的过程中，我也会继续提升自己的文字表达能力、画面审美能力和团队协作能力。

\par\qquad 恳请学生会给予我一次学习和实践的机会。感谢您的阅读！

\par\qquad 此致

\noindent 敬礼！

\par\hfill 申请人：小李

\par\hfill xx年xx月xx日
\end{aibox}

这次输出与提问之间有清楚对应。提问中提供了专业、经历、技能和加入后的打算，输出中也围绕这些内容展开。它没有突然增加获奖经历，也没有虚构大型项目。这说明，当学生向 AI 提供真实材料后，AI 更适合帮助学生进行组织和表达。

\subsubsection*{\textbf{• 成稿示例}}

AI 生成的申请信仍然是初稿。正式提交前，小李可以继续使用AI进行压缩和润色。比如有了更明确的兴趣可以改得更自然，继续提升自己的文字表达能力、画面审美能力和团队协作能力可以保留，但不宜再继续堆砌更多能力。下面是一版经过修改后的申请信。

\begin{aibox}
\noindent 尊敬的学生会负责人：

\par\qquad 您好！

\par\qquad 我是数字媒体技术专业的大二学生小李，现申请竞选学生会新媒体部干事。平时我比较关注校园活动的宣传与记录，也希望在学生会工作中进一步参与推文编辑、素材整理和短视频制作等实践。

\par\qquad 在以往的学习和活动中，我曾参与班级公众号图文排版，协助整理过三篇班级活动推文初稿，了解推文从素材整理到版面呈现的基本流程。校运动会期间，我也参与过活动照片拍摄和分类整理，体会到校园宣传不仅要记录现场，还要根据主题选择合适的图片和文字。除此之外，我平时会使用剪映制作短视频，也能用 Canva 完成简单海报设计。

\par\qquad 如果能够加入新媒体部，我希望从推文编辑、素材整理等基础工作做起，认真完成部门安排的任务，并在实践中提升图文表达、视频剪辑和团队协作能力。我也会注意内容的准确性和发布规范，努力为校园活动宣传贡献自己的力量。

\par\qquad 恳请给予我一次学习和实践的机会。感谢您的阅读！

\par\qquad 此致

\noindent 敬礼！

\par\hfill 申请人：小李

\par\hfill xx年xx月xx日
\end{aibox}

这版成稿比 AI 初稿更稳妥。它删去了部分重复表达，也把“从……做起”“认真完成”等态度放在具体任务中表达，而不是单独喊口号。申请信不是越华丽越好，而是要让阅读者相信：申请者确实了解这个岗位，也确实有一定基础和参与意愿。

前面的申请信示例已经把小李的经历连成了申请理由，但申请信能否使用，还取决于事实是否可靠、理由是否具体、语气是否得当。下面把案例中需要保留的审查知识集中说明。

\Needspace{9\baselineskip}
\begin{critiquebox}
申请信审查重点是“真实经历能否支撑申请理由”。AI 容易把普通经历写得过强，例如把整理推文初稿写成负责公众号运营，把基础剪辑写成熟练掌握视频制作。正式材料应保留事实边界，生成内容也需要人工核验其准确性和可靠性\footnotemark。

小李的申请信应围绕公众号排版、推文整理、运动会照片、剪映和 Canva 等具体材料展开，再说明加入后的打算。语气上避免两个极端：不要把自己写得过于完美，也不要写得没有基础。

\textbf{审查要点：}申请目的是否清楚，经历是否具体，顺序是否符合“申请目的—已有基础—后续打算”，语气是否真诚稳妥，是否写入无关个人信息或他人隐私。
\end{critiquebox}
\addtocounter{footnote}{-1}
\stepcounter{footnote}\footnotetext{参见：国家互联网信息办公室、国家发展和改革委员会、教育部等《生成式人工智能服务管理暂行办法》（2023年7月10日公布）第四条。访问日期：2026年7月18日。}



通过申请信，小李不仅展示了自己的经历，也表达了加入新媒体部后的工作方向。与简历相比，申请信更适合呈现完整的表达逻辑，也为竞选后的感谢表达提供了具体语境。


\subsection{感谢内容表达与关系维护}

感谢信的核心作用，是在他人提供指导、帮助或机会后进行真诚回应，并维护良好的后续关系。本小节围绕感谢信创作展开：先说明感谢信在老师指导、同学帮助、面试机会和组织录用等场景中的用途，再结合小李竞选成功后的情境，展示如何把感谢对象、具体原因和后续态度写清楚，并说明感谢信完成后需要从称呼、语气、具体性和隐私保护等方面进行审查。

感谢信在大学生活中也很常见。老师指导论文或比赛后，学生可以写感谢信；同学帮助准备材料后，也可以用简短感谢信表达心意；面试或竞选结束后，申请者可以向组织者表达感谢；节日或重要节点，也可以向老师、同学或合作伙伴表达问候。感谢信的作用不是堆砌礼貌话，而是表达对他人帮助的重视，并维护良好的后续关系。

在正式材料写作中，感谢信常常比简历和申请信更短，但并不意味着它可以随便写。很多 AI 生成的感谢信看起来很完整，却容易显得模板化。例如感谢您的帮助，我会继续努力，不辜负期望这样的句子虽然没有错，但如果没有具体事件，就会显得空泛。真正自然的感谢信，通常要包括感谢对象、具体帮助、自己的感受或后续态度。

\subsubsection*{\textbf{• 案例}}

小李通过了学生会新媒体部干事竞选。面试时，负责面试的学长学姐提醒她：新媒体部工作不仅要会做图文和视频，还要注意内容准确性，并且要和其他成员协作。竞选成功后，小李想给学长学姐写一封简短感谢信。她先进行了一个简单提问。

\begin{userbox}
帮小李写一封加入学生会后的感谢信。
\end{userbox}

\begin{aibox}
\noindent 学长学姐好：

\par\qquad 非常感谢大家给我这次加入学生会的机会。我会珍惜这次机会，认真学习，积极参与学生会工作，不断提升自己，努力完成各项任务。谢谢大家的认可和支持。

\par\qquad 再次感谢！

\par\hfill 小李

\par\hfill xx年xx月xx日
\end{aibox}

这封感谢信可以表达基本感谢，但它没有说明感谢谁，也没有说明具体感谢什么。由于提问中只说加入学生会后的感谢信，AI 只能生成通用表达。这样的内容适合临时参考，却不适合作为教材中的优秀示例。

\subsubsection*{\textbf{• 内容整理}}

写感谢信时，可以先想清楚三个小点：感谢对象、感谢原因和后续态度。感谢对象不同，称呼和语气也会不同；感谢原因不能只停留在谢谢，而要说明对方做了什么、提醒了什么或提供了什么帮助；后续态度则可以简要说明自己会如何继续学习、完成任务或保持联系。

小李重新整理了背景信息：她感谢的是负责面试的学长学姐；感谢原因是他们给了她加入新媒体部的机会，也在面试中提醒她注意推文内容准确性和团队协作；她后续希望认真参与推文编辑、素材整理和短视频制作。整理清楚后，她再次提问。

\begin{userbox}
请帮小李写一封加入学生会新媒体部后的感谢信，收信对象是负责面试的学长学姐。

背景是：小李通过了学生会新媒体部干事竞选。面试时，学长学姐提醒她，新媒体部工作要注意推文内容的准确性，也要重视团队协作。小李希望表达感谢，并说明自己之后会认真参与推文编辑、素材整理和短视频制作。

要求语气真诚自然，不要太官方，150 字左右。
\end{userbox}

\begin{aibox}
\noindent 学长学姐好：

\par\qquad 非常感谢你们在面试中给予我的建议，也感谢新媒体部给我这次加入和学习的机会。面试时你们提到推文内容要注意准确性，也要重视团队协作，这让我对新媒体部的工作有了更清楚的认识。之后我会认真参与推文编辑、素材整理和短视频制作，主动向大家学习，努力把分配到的任务做好。

\par\qquad 再次感谢你们的认可和帮助，祝学习和工作顺利！

\par\hfill 小李

\par\hfill xx年xx月xx日
\end{aibox}

这一次生成的感谢信更自然，也更有针对性。它没有简单套用感谢认可这样的表达，而是提到了面试中的具体建议。这样对方能够看出，小李确实认真听取了建议，也愿意在后续工作中落实。

\subsubsection*{\textbf{• 成稿示例}}

感谢信通常不需要太长。若是发在微信或 QQ 中，可以比正式书信更简短；若是发给老师或正式组织负责人，则应更加稳重。小李要发给负责面试的学长学姐，因此可以保留亲切自然的语气。

\begin{aibox}
\noindent 学长学姐好：

\par\qquad 谢谢你们在面试中给我的建议，也谢谢新媒体部给我这次加入和学习的机会。你们提到推文内容要准确、部门工作要注意协作，这让我对之后的工作有了更清楚的认识。接下来我会认真参与推文编辑、素材整理和短视频制作，主动向大家学习，把分配到的任务做好。

\par\qquad 再次感谢你们的认可和帮助，祝学习和工作顺利！

\par\hfill 小李

\par\hfill xx年xx月xx日
\end{aibox}

这版成稿去掉了部分重复表达，比 AI 初稿更适合直接发送。感谢信的关键不是写得多正式，而是让对方感到真诚。若小李要把这封信发给老师，可以把学长学姐好改为老师您好，并适当提高语气的正式程度。若发给同学，则可以更轻松一些。

前面的感谢信示例已经能够表达谢意，但感谢信不能只追求礼貌和完整。小李在发送前，还需要根据对象关系、感谢原因和公开边界进行检查。

\Needspace{9\baselineskip}
\begin{critiquebox}
感谢信审查不看辞藻是否华丽，而看对象、原因和分寸是否合适。写给老师、学长学姐和组织负责人，称呼和语气不能完全一样；感谢原因也不能只写“感谢帮助”，应点明对方提供的建议、机会或具体支持。

AI 常会生成固定套话，如“在今后的学习和工作中不负期望”。这类表达可以保留一两句，但不能堆砌。感谢信还应避免写入他人手机号、私人行程、未公开评价等信息，使用 AI 起草时也要保持最小必要意识\footnotemark。

\textbf{审查要点：}称呼是否匹配对象，感谢原因是否具体，语气是否自然，结尾是否适合关系远近，是否泄露自己或他人的非必要信息。
\end{critiquebox}
\addtocounter{footnote}{-1}
\stepcounter{footnote}\footnotetext{参见：全国人民代表大会常务委员会《中华人民共和国个人信息保护法》（2021年8月20日公布，2021年11月1日起施行）第五条、第六条。访问日期：2026年7月18日。}



通过感谢信，小李完成了竞选过程中的最后一类正式材料。简历帮助她展示基础，申请信帮助她说明理由，感谢信帮助她表达态度和维护关系。三个文本连在一起，形成了一个完整的应用场景，也展示了 AI 在正式材料写作中的辅助价值。

\subsection*{小结}
本节围绕简历材料、申请陈述和感谢表达展开，重点说明如何把个人经历转化为真实、清楚、得体的书面材料。简历写作要突出经历与岗位、项目或申请目标之间的关联，申请陈述要说明动机、基础和后续计划，感谢表达则要兼顾礼貌、具体和分寸。使用 AI 时，关键不是让模型替代个人经历，而是帮助整理信息、优化结构和检查表达。生成后仍要核对时间、身份、成果和语气，避免夸大经历、虚构事实或使用过度模板化的语言，使材料更具体、更可信。
\bigskip

\subsection*{课后习题}
\begin{enumerate}
\setlength{\itemsep}{6pt}

\item 选择一个真实的学习、竞选、实习或项目报名场景，整理一份可提供给 AI 的个人经历素材清单，至少包括时间、身份、任务、过程和结果。

\item 根据上述素材，设计一段用于生成简历条目或申请陈述的 AI 提问，要求写明使用对象、申请目标、语气要求和输出格式。

\item 对比简历、申请信和感谢信三类正式材料，说明它们在写作目的、内容重点和语气分寸上的差异。

\item 检查一段 AI 生成的正式材料时，应重点核对哪些事实信息？请从经历真实性、成果表述、时间地点和身份关系等方面说明。

\item 结合正式材料写作场景，说明为什么 AI 输出不能直接当作最终材料提交，并谈谈使用者应承担哪些审查责任。

\end{enumerate}

\newpage

\section{办公材料整理与文档写作}

\begin{sectionkeypoints}
\item 理解办公材料整理与文档写作的基本任务，认识活动材料、会议记录、反馈意见和经费清单等原始资料的整理价值。
\item 掌握使用 AI 提取和归纳办公信息的基本方法，能够围绕时间、地点、人数、名单、金额、任务和责任人等要素进行核对。
\item 了解总结、方案、通知、请示、函和会议纪要等常见文种的用途差异，理解已完成工作汇报和下一步工作安排的表达重点。
\item 熟悉 AI 在办公写作中的辅助流程，能够利用其完成材料提取、结构梳理、初稿起草和语言润色。
\item 认识办公写作的责任边界，明确涉及数据、经费、名单和正式发布内容时必须回到原始材料进行人工确认。
\end{sectionkeypoints}

办公写作场景关注的是学生在学习、社团、学生组织和项目实践中，如何把已经发生的事情、已经收集的材料和后续需要推进的任务，整理成可以提交、汇报、通知或执行的正式文本。与创意写作不同，办公写作不追求语言新奇，也不是单纯让文章看起来完整，而是要求内容真实、结构清楚、事项明确、责任可追踪。一次活动结束后，学生往往需要整理材料、归纳问题、撰写总结、形成方案、发布通知、准备请示或撰写会议纪要，这些工作共同构成了较完整的办公写作链条。

从实际使用场景看，办公写作通常具有连续性。活动签到表、报名表、照片说明和现场反馈，可以先被整理成事实清单；事实清单可以转化为活动总结；总结中发现的问题，可以进一步形成改进方案；如果下一次活动需要场地、经费或设备，还可能继续生成请示、通知或函件。因此，办公写作不是一个孤立的写作动作，而是围绕具体事务不断推进的文本处理过程。

本节以一次校园 AI 创意作品展示活动结束后的后续工作为线索展开。活动结束以后，小李所在的新媒体部需要先整理活动资料，再撰写活动总结和改进方案，并根据后续工作需要起草通知、请示、函件和会议纪要。在这一过程中，AI 可以帮助小李提取关键信息、搭建文档结构，并对语言进行初步调整；但活动人数、作品数量、经费金额、责任分工和提交时间等关键内容，仍然必须回到原始材料中核对。

围绕这个案例，本节依次介绍材料提取与要点归纳、总结与方案撰写、公文起草与纪要整理。通过这一流程，学生可以理解 AI 在办公写作中的正确定位：它可以辅助整理、生成和润色，但不能替代事实核对、责任判断和人工审查。


\subsection{办公材料提取与要点归纳}

在学习和工作中，办公写作往往不是从空白页面开始的。很多时候，学生面对的是一堆已经发生过的事情、一批零散的记录和若干需要后续处理的材料。例如一次校园活动结束后，可能会有报名表、签到表、活动照片、主持稿、现场记录、经费清单、反馈问卷、工作人员聊天记录和会议讨论内容。如果直接开始写总结或方案，很容易遗漏信息，也容易把不同来源的内容混在一起。

材料提取与要点归纳的作用，就是把零散材料先整理成可以继续使用的信息。它可以用于会议记录整理、活动材料汇总、课程小组分工、调研资料初步分析、项目过程记录、实习日志归纳等场景。对学生来说，这一步看似普通，却直接影响后续写作质量。如果前期材料整理不清楚，后面写出来的总结、方案和纪要就容易空泛、不准确，甚至出现事实错误。

\subsubsection*{\textbf{• 案例}}

小李成功加入学生会新媒体部后，学生会举办了一场校园 AI 创意作品展示活动，活动结束后，老师要求新媒体部协助整理活动材料，并在此基础上完成活动总结、改进方案、公文起草和会议纪要等后续工作。小李需要先处理的不是直接写成文档，而是把报名表、签到表、活动照片、现场记录、经费清单、反馈问卷和会议讨论内容等资料进行初步整理。接下来，本小节将以这一活动为例，说明如何借助 AI 提取关键信息、归纳材料要点，并从信息来源、事实判断和数据核对三个方面完成后续审查。
小李一开始使用了一个很短的提问。

\begin{userbox}
下面是校园 AI 创意作品展示活动的一些材料，请帮我整理一下。

材料：活动在周五下午举行，地点是大学生活动中心一楼展厅。现场有很多同学来参观，有 AI 绘画、AI 海报、AI 故事和 AI 配音作品。老师说整体效果不错，但现场引导有点乱。工作人员说照片还没有整理完，优秀作品名单也还没有最后确认。
\end{userbox}

\begin{aibox}
本次校园 AI 创意作品展示活动整体效果较好，活动内容丰富，包含 AI 绘画、AI 海报、AI 故事和 AI 配音等作品形式。活动地点为大学生活动中心一楼展厅，时间为周五下午。现场同学参与积极，老师对活动效果给予肯定。

后续需要改进的地方主要包括现场引导秩序、照片整理和优秀作品名单确认。建议工作人员尽快完成活动照片归档，确认优秀作品名单，并在后续活动中提前安排引导人员，优化现场参观流程。
\end{aibox}

这个回答能够大致概括材料，但它仍然比较粗。它没有区分已经确定的信息和需要核实的信息，也没有把任务整理成清单。例如周五下午并不是规范时间，现场有很多同学来参观也不是准确人数。若直接把这段内容用于活动总结，就可能显得不够正式。

\subsubsection*{\textbf{• 内容整理}}

使用 AI 整理材料时，学生要先说明整理目的，而不是只要求 AI 概括一下。材料整理的目标，是把零散记录转化为后续写总结、方案、新闻稿或纪要时可以直接使用的信息。因此，在向 AI 提交材料前，需要先说明这批材料要用于什么写作任务，并要求 AI 按照固定栏目进行归类。

已有材料也要尽量完整地列出。材料可以包括活动名称、时间地点、参与情况、现场情况、反馈意见、未完成事项和后续安排等内容。如果某些信息不完整，也应如实写出，而不是让 AI 自动补齐。对于没有准确数据的内容，要明确要求 AI 标注待核实，避免把推测内容写成事实。

整理格式也要提前设定。学生可以要求 AI 按基本事实、活动成果、存在问题、后续任务等类别输出，也可以要求使用表格呈现。这样整理后的内容更清楚，后续撰写活动总结、改进方案或相关公文时也更方便调用。小李随后把校园 AI 创意作品展示活动的材料补充得更具体，并要求 AI 按照表格进行整理。

\begin{userbox}
请把下面这组校园 AI 创意作品展示活动材料整理成办公写作可用的材料清单。要求分为基本事实、活动成果、存在问题、后续任务四类。凡是材料中没有准确数据的地方，请标注待核实，不要自行编造。

已知材料如下：
1. 活动名称：校园 AI 创意作品展示活动。
2. 活动时间：周五下午，具体日期需要再查活动通知。
3. 活动地点：大学生活动中心一楼展厅。
4. 展示内容：AI 绘画、AI 海报、AI 故事、AI 配音作品。
5. 现场情况：参观同学较多，但没有准确统计人数。
6. 老师反馈：整体效果较好，作品形式比较丰富，但现场引导秩序需要改进。
7. 工作人员反馈：照片还没有整理完，优秀作品名单还没有最终确认。
8. 后续工作：需要整理活动照片、确认优秀作品名单、撰写活动总结、制作获奖证书、准备下一次活动改进方案。
\end{userbox}

\begin{aibox}
\textbf{一、基本事实}

活动名称：校园 AI 创意作品展示活动。

活动时间：周五下午，具体日期待核实。

活动地点：大学生活动中心一楼展厅。

展示内容：AI 绘画、AI 海报、AI 故事、AI 配音作品。

\textbf{二、活动成果}

1. 活动展示形式较丰富，包含图像、海报、文字故事和配音等多种 AI 创作成果。

2. 现场有较多同学参观，说明活动具有一定吸引力，但具体参观人数待核实。

3. 老师认为活动整体效果较好，作品形式较丰富。

\textbf{三、存在问题}

1. 现场引导秩序需要改进。

2. 活动照片尚未完成整理。

3. 优秀作品名单尚未最终确认。

4. 活动时间中的具体日期需要根据原始通知核实。

\textbf{四、后续任务}

1. 查询活动通知，补全准确日期。

2. 整理并归档活动照片。

3. 确认优秀作品名单。

4. 撰写活动总结。

5. 制作获奖证书。

6. 准备下一次活动改进方案。
\end{aibox}

第二次输出更加适合办公写作。它把材料分成了事实、成果、问题和任务，也把具体日期、参观人数等不确定信息标为待核实。这样的整理结果不能直接作为最终总结，但它可以作为后续写总结、方案和纪要的材料基础。

\subsubsection*{\textbf{• 成稿示例}}

小李在 AI 整理结果的基础上，又回到原始材料中核对了活动通知和签到表。她确认活动时间为 xx年xx月xx日 14:30 至 16:30；签到表显示现场签到人数为 86 人；展示作品共 32 件，其中 AI 绘画作品 12 件、AI 海报作品 8 件、AI 故事作品 7 件、AI 配音作品 5 件。根据这些核实后的信息，她可以形成一份更清楚的材料汇总表。核实后的活动材料整理结果见表\ref{tab:ch4_05_ai_activity_material}。

\begin{formalLongTable}
\begin{longtable}{p{0.22\textwidth}p{0.68\textwidth}}
\caption[校园 AI 创意作品展示活动材料整理表]{\chaptertablelabel\quad 校园 AI 创意作品展示活动材料整理表}
\label{tab:ch4_05_ai_activity_material}\\
\toprule
\textbf{类别} & \textbf{整理结果} \\
\midrule
\endfirsthead

\multicolumn{2}{c}{\chaptertablelabel\quad 校园 AI 创意作品展示活动材料整理表（续）} \\
\toprule
\textbf{类别} & \textbf{整理结果} \\
\midrule
\endhead

\midrule
\multicolumn{2}{r}{续下页} \\
\endfoot

\bottomrule
\endlastfoot

活动名称 & 校园 AI 创意作品展示活动 \\
活动时间 & xx年xx月xx日 14:30 至 16:30 \\
活动地点 & 大学生活动中心一楼展厅 \\
参与情况 & 现场签到人数 86 人，参与展示作品 32 件 \\
作品类型 & AI 绘画 12 件、AI 海报 8 件、AI 故事 7 件、AI 配音 5 件 \\
活动亮点 & 展示形式较丰富，作品覆盖图像、文字和声音等多种表达形式，现场参观积极性较高 \\
存在问题 & 现场引导标识不够明显，照片归档不及时，优秀作品名单确认较晚 \\
后续任务 & 完成照片整理、确认优秀作品名单、制作证书、撰写总结、提出下一次活动改进方案 \\

\end{longtable}
\end{formalLongTable}

这张表的作用，是把零散材料变成后续写作可以直接调用的信息。相比原始材料，它更准确、更有层次，也便于老师和同学快速查看。

前面的材料整理示例把零散信息转成了事实清单，但事实清单并不等于最终事实。小李继续使用这些材料前，需要先区分已确认信息、待核实信息和主观评价。

\Needspace{9\baselineskip}
\begin{critiquebox}
材料提取审查的核心是分清“事实、评价、待核实”。活动时间要查通知，人数要查签到表，经费要查票据，获奖名单要查最终确认文件；“活动效果较好”属于评价，不能和时间、地点、人数这类事实混写。

无法确认的信息不要让 AI 补齐，应标注“待核实”。人数、作品数量、经费金额等数据尤其要回到原始材料中复核，保证真实、准确、完整和及时\footnotemark。

\textbf{审查要点：}是否遗漏时间地点，人数是否有签到表支撑，成果数量是否按清单统计，个别反馈是否被扩大成整体结论，后续任务是否写清责任人和时间节点。
\end{critiquebox}
\addtocounter{footnote}{-1}
\stepcounter{footnote}\footnotetext{参见：全国人民代表大会常务委员会《中华人民共和国统计法》（2024年9月13日修订通过并公布）第一条。访问日期：2026年7月18日。}



完成材料提取后，小李就有了比较可靠的写作基础。接下来，她需要用这些材料撰写活动总结，并进一步提出下一次活动的改进方案。


\subsection{活动总结撰写与方案设计}

活动总结和活动方案是校园办公写作中常见的两类事务性文本，但二者的写作方向不同。活动总结面向已经完成的工作，重点是对活动过程、参与情况、实际成效和存在问题进行回顾与分析，其核心不是简单复述活动经过，而是从事实材料中提炼经验和不足。活动方案面向尚未开展或需要继续推进的工作，重点是明确目标、任务、人员、时间、资源和保障措施，其核心不是提出笼统想法，而是把后续工作安排转化为可以执行的行动计划。

学生在社团活动、班级建设、志愿服务、竞赛组织、课程项目和实习实践中，经常需要写总结和方案。例如一次活动结束后，要写活动总结报给老师；发现活动中存在问题后，要写改进方案；下一次活动开始前，要写执行方案；需要申请场地或经费时，还要把方案作为附件提交。AI 可以帮助学生搭建结构、整理语言、补全过渡，但它不能替代真实材料。

本小节围绕总结与方案写作展开，重点说明如何把已经整理好的活动材料转化为正式办公文本。以校园 AI 创意作品展示活动的后续工作为例，学生需要先依据签到表、作品清单、现场反馈等材料形成活动总结，再根据总结中暴露的问题提出下一次活动的改进方案。通过这一过程，可以看到：总结强调事实支撑和问题归纳，方案强调任务分工和时间节点，二者共同服务于活动工作的复盘与改进。

\subsubsection*{\textbf{• 活动总结案例}}

小李已经完成了校园 AI 创意作品展示活动的材料整理。接下来，她需要把材料转化为两类文本：一类是活动总结，用来向老师说明活动开展情况；另一类是下一次活动的改进方案，用来解决现场引导、照片归档和名单确认等问题。小李第一次让 AI 写总结时，只给了一个很短的要求。

\begin{userbox}
帮我写一份校园 AI 创意作品展示活动总结。
\end{userbox}

\begin{aibox}
校园 AI 创意作品展示活动总结

本次校园 AI 创意作品展示活动顺利开展，活动内容丰富，形式多样，充分展示了同学们在人工智能创意应用方面的学习成果。活动过程中，同学们积极参与，现场氛围良好，作品展示效果较好。

通过本次活动，同学们进一步了解了 AI 在绘画、海报设计、故事创作和配音等方面的应用，也提高了对人工智能技术的兴趣。活动不仅丰富了校园文化生活，也为同学们提供了展示创意和交流学习的平台。

当然，活动中也存在一些不足，如现场秩序维护还需加强，宣传和组织细节仍有改进空间。今后，我们将进一步完善活动流程，提升组织水平，努力举办更多高质量校园活动。
\end{aibox}

这份总结比较通用，适合做参考框架，但不适合直接提交。因为它没有具体时间、地点、人数和作品数量，也没有说明问题来自哪里。更重要的是，它把活动意义写得比较大，却没有呈现活动过程中的真实细节。

\subsubsection*{\textbf{• 活动总结内容整理}}

使用 AI 生成活动总结前，学生需要先把已整理的材料进一步核实清楚。活动总结属于正式材料，不能只给 AI 一些零散印象，而应提供活动名称、具体时间、地点、参与人数、作品数量、活动内容、反馈意见和存在问题等事实信息。对于已经核实的数据，可以直接写入材料；对于没有核实的信息，则不要让 AI 自行补充。

在向 AI 提出写作要求时，还要说明总结的结构和表达边界。可以要求 AI 按活动概况、主要内容、活动成效、存在问题、改进方向几个部分展开，并使用正式、具体的语言。活动成效要有数据和反馈支撑，存在问题要来自真实材料，改进方向要与问题对应，不能泛泛而谈。小李在前面已经完成了材料整理和数据核实，于是根据已核实材料重新向 AI 提问。

\begin{userbox}
请根据以下已核实材料，帮我写一份校园 AI 创意作品展示活动总结。要求语言正式、具体，不要夸大，不要编造材料。结构包括活动概况、主要内容、活动成效、存在问题、改进方向五个部分。

已核实材料：活动名称为校园 AI 创意作品展示活动；活动时间为 xx年xx月xx日 14:30 至 16:30；地点为大学生活动中心一楼展厅；现场签到人数 86 人；展示作品共 32 件，包括 AI 绘画 12 件、AI 海报 8 件、AI 故事 7 件、AI 配音 5 件。老师反馈整体效果较好，作品形式较丰富，但现场引导秩序需要改进。工作人员反馈活动照片整理不够及时，优秀作品名单确认较晚。
\end{userbox}

\begin{aibox}
校园 AI 创意作品展示活动总结

\textbf{一、活动概况}

xx年xx月xx日 14:30 至 16:30，校园 AI 创意作品展示活动在大学生活动中心一楼展厅举行。本次活动面向全校学生展示 AI 创意作品，现场签到人数 86 人，共展出作品 32 件。

\textbf{二、主要内容}

本次活动展示内容包括 AI 绘画、AI 海报、AI 故事和 AI 配音四类作品。其中，AI 绘画作品 12 件，AI 海报作品 8 件，AI 故事作品 7 件，AI 配音作品 5 件。活动现场通过作品陈列和现场讲解等方式，展示了同学们对人工智能创意应用的探索。

\textbf{三、活动成效}

从现场情况看，本次活动参与度较高，展示内容较丰富，能够让参观同学直观了解 AI 在图像、文字和声音等不同创作形式中的应用。老师反馈认为，本次活动整体效果较好，作品形式比较多样，具有一定的展示和交流意义。

\textbf{四、存在问题}

活动中也存在一些需要改进的地方。现场引导秩序还不够清晰，参观路线和作品分区提示需要进一步优化；活动照片整理不够及时，影响了后续宣传材料制作；优秀作品名单确认较晚，也使证书制作和结果公布受到一定影响。

\textbf{五、改进方向}

后续开展类似活动时，应提前设置现场引导标识，明确参观路线和作品分区；安排专人负责照片收集、筛选和归档；在活动前明确优秀作品评选流程和名单确认时间，减少活动结束后的材料整理压力。
\end{aibox}

第二次总结更适合办公写作。它使用了已核实数据，结构也比较清楚。可以看到，AI 的主要作用不是凭空写总结，而是把材料转化为规范文本。材料越清楚，生成结果越可用。

\subsubsection*{\textbf{• 活动方案案例}}

活动方案与活动总结不同。活动总结回顾已经发生的事情，方案安排即将开展的事情。小李在总结中发现，本次活动存在现场引导不够清晰、照片整理不及时、优秀作品名单确认较晚等问题。为了下一次活动更顺利，她需要写一份改进方案。


小李先进行了一个简单提问。

\begin{userbox}
帮我写一份下一次校园 AI 活动的改进方案。
\end{userbox}

\begin{aibox}
下一次校园 AI 活动改进方案

为进一步提升校园 AI 活动质量，增强活动吸引力和组织效果，下一次活动应在宣传、组织、现场管理和后期总结等方面进行改进。活动前要加强宣传推广，扩大同学参与度；活动中要做好现场秩序维护，安排工作人员进行引导；活动后要及时整理资料，总结经验，不断提高活动水平。

通过以上措施，下一次校园 AI 活动将更加规范有序，内容更加丰富，参与效果更加明显。
\end{aibox}

这个方案方向正确，但仍然比较泛。它没有回应本次活动中具体出现的问题，也没有写明人员分工和时间安排。小李再次根据活动总结中的问题进行提问。

\subsubsection*{\textbf{• 活动方案内容准备}}

使用 AI 策划活动方案时，学生不能只要求 AI 写一份方案，而要先说明方案要解决什么问题。活动方案不是把目标写得好听，而是把问题转化为可以执行的安排。因此，在向 AI 提出要求时，应先交代活动复盘中发现的主要问题，例如现场组织、材料整理、人员分工、时间衔接、结果确认等方面的不足，再要求 AI 围绕这些问题提出改进措施。

同时，要向 AI 说明方案的基本背景和输出结构。背景包括活动由谁协助组织、下一次活动大致在哪里举办、包含哪些主要环节；输出结构可以包括改进目标、具体措施、人员分工、时间安排和检查方式。这样生成的方案才不会停留在加强管理、提高质量等笼统表述上，而能形成谁负责、什么时候完成、如何检查的具体安排。小李根据活动复盘情况，围绕下一次活动需要解决的问题，向 AI 提出了更明确的方案生成要求。


\begin{userbox}
请根据校园 AI 创意作品展示活动复盘情况，写一份下一次活动改进方案。要求重点解决三个问题：现场引导秩序不够清晰、活动照片整理不及时、优秀作品名单确认较晚。方案要包括改进目标、具体措施、人员分工、时间安排和检查方式。语言正式、可执行，不要空泛。

补充背景：活动由学生会新媒体部协助组织。下一次活动仍在校园内举办，预计包含作品展示、现场参观和优秀作品评选三个环节。
\end{userbox}

\begin{aibox}
下一次校园 AI 创意活动改进方案

\textbf{一、改进目标}

为提升校园 AI 创意活动的组织质量，针对上次活动中现场引导不够清晰、照片整理不及时、优秀作品名单确认较晚等问题，下一次活动将重点优化现场管理、素材归档和评选流程，确保活动组织更加有序，后续材料整理更加及时。

\textbf{二、具体措施}

1. 优化现场引导。活动前制作场地分区图和参观路线提示，在入口、作品展示区和出口处设置明显标识，并安排工作人员进行现场引导。

2. 完善素材整理。活动当天指定专人负责照片拍摄和素材收集，活动结束后当天完成照片初步筛选，次日完成分类归档。

3. 提前明确评选流程。活动前确定优秀作品评选标准、评委人员和名单确认时间，活动结束后两日内完成优秀作品名单复核。

\textbf{三、人员分工}

现场引导组负责参观路线设计和秩序维护；素材整理组负责照片拍摄、筛选和归档；评选协调组负责优秀作品名单汇总和确认；文字宣传组负责活动总结和后续推文初稿。

\textbf{四、时间安排}

活动前一周完成现场分区和评选流程确认；活动前两天完成引导标识制作；活动当天完成现场引导和素材采集；活动结束后一天内完成照片初筛；活动结束后两天内完成优秀作品名单确认；活动结束后三天内完成总结和宣传材料初稿。

\textbf{五、检查方式}

活动负责人在活动前检查引导标识、人员分工表和评选流程表；活动结束后检查照片归档情况、优秀作品名单确认情况和总结材料完成情况，确保每项任务有记录、有责任人、有完成时间。
\end{aibox}

这份方案比第一次更可执行。它不是泛泛地说加强管理，而是把问题转化为具体任务，并安排了时间节点和人员分工。学生使用 AI 写方案时，应该让 AI 围绕真实问题展开，而不是只让它写一篇看起来完整的方案。

前面的总结和方案示例已经把活动材料转化为正式文本，但这两类材料不能只看格式是否完整。小李还需要判断总结是否基于事实、方案是否能够执行。

\begin{critiquebox}
总结和方案的审查标准不同。总结重在事实支撑和问题归纳，不能只写“活动圆满成功”；方案重在可执行性，不能只写“加强组织”“提升体验”这类空话。

小李修改时，应让总结中的每个成效都有材料来源，让方案中的每项措施对应真实问题。若任务没有责任人、完成时间和反馈方式，方案就只是愿望清单。

\textbf{审查要点：}总结是否有数据和材料来源，问题是否对应改进方向；方案是否有目标、任务、分工和时间节点，措施是否能够落地。
\end{critiquebox}


\Needspace{9\baselineskip}
正式办公材料尤其要重视实事求是和可操作性。《党政机关公文处理工作条例》要求公文起草一切从实际出发，分析问题实事求是，所提政策措施和办法切实可行；同时要求内容简洁、主题突出、结构严谨、表述准确\footnote{参见：中共中央办公厅、国务院办公厅《党政机关公文处理工作条例》（中办发〔2012〕14号，2012年4月16日发布，2012年7月1日起施行）第十九条。访问日期：2026年7月18日。}。学生写活动总结和改进方案虽然不是党政机关公文，但也可以借鉴这种审查原则：总结要基于真实材料，方案要能落实，文字要清楚精练。表\ref{tab:ch4_07_summary_plan_compare}用于说明总结和方案的区别。

\begin{formalLongTable}
\begin{longtable}{p{0.22\textwidth}p{0.33\textwidth}p{0.35\textwidth}}
\caption[活动总结与活动方案比较表]{\chaptertablelabel\quad 活动总结与活动方案比较表}
\label{tab:ch4_07_summary_plan_compare}\\
\toprule
\textbf{比较项目} & \textbf{活动总结} & \textbf{活动方案} \\
\midrule
\endfirsthead

\multicolumn{3}{c}{\chaptertablelabel\quad 活动总结与活动方案比较表（续）} \\
\toprule
\textbf{比较项目} & \textbf{活动总结} & \textbf{活动方案} \\
\midrule
\endhead

\midrule
\multicolumn{3}{r}{续下页} \\
\endfoot

\bottomrule
\endlastfoot

写作对象 & 已经完成的活动 & 即将开展或需要改进的活动 \\
核心问题 & 做了什么，效果如何，有什么不足 & 为什么做，做什么，谁来做，什么时候完成 \\
材料基础 & 活动记录、数据、反馈、照片、签到表 & 总结中发现的问题、后续目标、资源条件 \\
语言重点 & 客观回顾，适度评价 & 明确安排，便于执行 \\
常见问题 & 数据不准确，评价空泛，问题不具体 & 措施笼统，分工不清，时间节点缺失 \\

\end{longtable}
\end{formalLongTable}

无论写总结还是方案，都不能把 AI 的输出直接当作最终稿。活动总结中的人数、作品数量、时间地点必须核对；活动方案中的任务、负责人和完成时间必须与实际条件匹配。AI 可以帮助学生把材料组织成文本，但学生要负责判断文本是否真实、准确、可执行。


通过总结和方案，小李完成了活动结束后的第二类办公写作任务。总结让老师了解活动结果，方案让下一次活动有明确改进方向。但在学生组织工作中，除了总结和方案，还经常需要通知、请示、函和会议纪要等公文类材料。


\subsection{常用公文起草与格式审查}

公文和会议纪要是校园办公写作中较为正式的事务性文本。它们的重点不只是语言正式，更在于文种准确、事项清楚、结构规范、责任明确。通知、请示、函和会议纪要虽然都属于办公材料，但用途并不相同。通知主要用于告知事项和提出要求，请示主要用于向上级请求批准，函主要用于平级或不相隶属单位之间联系事项，会议纪要主要用于记录会议讨论内容和形成的决定。文种不同，写作结构、语气表达和结尾方式也不同，不能简单套用同一个模板。

在学生学习、社团和学生组织工作中，这类材料经常出现在活动推进和后续整理环节。例如，活动结束后需要发布材料补交通知，需要向学院提交经费请示，需要向其他部门发送协作函，也可能需要整理复盘会议纪要。与普通说明文字相比，这些材料更强调事实依据和执行要求。时间、地点、对象、材料名称、截止时间、联系人、经费金额和责任人等信息，都必须来自真实材料，不能由 AI 自行补写。

本小节围绕公文起草与会议纪要整理展开，重点说明 AI 如何辅助完成更正式的办公材料。具体来说，本小节先介绍通知、请示、函和会议纪要的不同用途，再通过活动后续通知、经费请示、材料报送函和复盘会议纪要进行示范，最后从文种、格式、事实、语气和责任五个方面说明打磨与审查要求。

\subsubsection*{\textbf{• 案例}}

在校园 AI 创意作品展示活动结束后，小李所在的新媒体部还需要完成一组后续工作：面向参与同学发布优秀作品名单确认通知，向指导老师提交证书制作和材料打印的经费请示，向学院团委报送活动成果，并整理复盘会议纪要、明确后续任务。小李可以使用 AI 帮助起草这些材料，但必须人工核对格式和事实。

\subsubsection*{\textbf{• 通知的创作与审查}}

通知用于向特定对象传达事项安排，重点是让接收者知道需要做什么、什么时候完成、通过什么方式反馈。小李所在部门需要在校园 AI 创意作品展示活动结束后，提醒入选优秀作品初步名单的同学核对个人信息。通知不能只写活动已经结束，而要说明通知对象、核对内容、截止时间、反馈方式和逾期处理办法。

小李先用一个简单提问测试 AI。

\begin{userbox}
帮我写一个校园 AI 创意作品展示活动后续通知。
\end{userbox}

\begin{aibox}

\begin{center}
\textbf{关于校园 AI 创意作品展示活动后续工作的通知}
\end{center}


\noindent 各位同学：

\par\qquad 校园 AI 创意作品展示活动已经顺利结束。为做好后续工作，请参与活动的同学关注作品评选、证书领取和材料整理等相关安排。请大家按照要求及时配合，确保后续工作顺利进行。

\par\qquad 特此通知。

\par\hfill 学生会新媒体部

\par\hfill xx年xx月xx日
\end{aibox}

这个通知只是一个空框架。它没有说明谁需要做什么，也没有截止时间和联系方式。若直接发布，同学们不知道该确认名单、提交信息还是等待证书。因此，小李需要把具体事项提供给 AI。

\begin{userbox}
请帮我起草一则通知，用于校园 AI 创意作品展示活动结束后，提醒入选优秀作品名单的同学核对个人信息。要求格式正式，事项清楚。

具体信息如下：通知对象为入选优秀作品初步名单的同学；需要核对姓名、学院班级、作品名称和联系方式；核对截止时间为 xx年xx月xx日 18:00；核对方式是发送信息给学生会新媒体部邮箱 aiworks@example.com；逾期未反馈的，默认信息无误。落款为学生会新媒体部，日期为 xx年xx月xx日。
\end{userbox}

\begin{aibox}
\begin{center}
\textbf{关于核对校园 AI 创意作品展示活动优秀作品信息的通知}
\end{center}

\noindent 各位同学：

\par\qquad 校园 AI 创意作品展示活动优秀作品初步名单已完成整理。为保证后续证书制作和结果公布信息准确，请入选优秀作品初步名单的同学及时核对个人信息。

\par\qquad 请相关同学核对姓名、学院班级、作品名称和联系方式，并于 xx年xx月xx日 18:00 前将核对结果发送至学生会新媒体部邮箱 aiworks@example.com。逾期未反馈的，默认相关信息无误。

\par\qquad 请各位同学认真核对，感谢配合。

\par\hfill 学生会新媒体部

\par\hfill xx年xx月xx日
\end{aibox}

这一次生成的通知更符合实际使用。它明确了通知对象、核对内容、截止时间、反馈方式和逾期处理方式。通知类公文的关键是让接收者知道自己要做什么，而不是只写活动已经结束。

\subsubsection*{\textbf{• 请示的创作与审查}}

请示用于向上级请求批准，通常一事一请。小李所在部门需要制作优秀作品证书，并打印部分展示材料。因为涉及经费支出，需要向指导老师或学院提交请示。请示不能只写我们想申请经费，而要说明原因、用途、金额和请求事项。

\begin{userbox}
帮我写一份申请活动经费的请示。
\end{userbox}

\begin{aibox}

\begin{center}
\textbf{关于申请活动经费的请示}
\end{center}
\noindent 尊敬的老师：

\par\qquad 为保障校园 AI 创意作品展示活动顺利开展，现申请一定活动经费用于相关物资采购和材料制作。恳请老师审批。

\par\qquad 妥否，请批示。

\par\hfill 学生会新媒体部

\par\hfill xx年xx月xx日
\end{aibox}

这份请示不够可用。它没有说明活动已经结束，费用用于什么，也没有具体金额。若提交给老师，老师无法判断是否合理。小李补充了具体信息后再次提问。

\begin{userbox}
请帮我起草一份经费请示。背景是校园 AI 创意作品展示活动已经结束，现需要制作优秀作品证书 20 份，并打印活动成果汇编材料 10 份。预计证书制作费用 100 元，材料打印费用 80 元，合计 180 元。请示对象是学院团委老师，落款为学生会新媒体部，日期为 xx年xx月xx日。要求文种正确、语气正式，结尾使用请示类结尾。
\end{userbox}

\begin{aibox}
\begin{center}
\textbf{关于申请校园 AI 创意作品展示活动后续材料制作经费的请示}
\end{center}

\noindent 学院团委老师：

\par\qquad 校园 AI 创意作品展示活动已顺利结束。为做好活动后续成果整理与优秀作品表彰工作，现需制作优秀作品证书 20 份，并打印活动成果汇编材料 10 份。

\par\qquad 经初步核算，证书制作费用预计 100 元，材料打印费用预计 80 元，合计 180 元。该项费用主要用于活动后续材料制作和成果归档。

\par\qquad 现申请活动后续材料制作经费 180 元，请予审批。

\par\qquad 妥否，请批示。

\par\hfill 学生会新媒体部

\par\hfill xx年xx月xx日
\end{aibox}

这份请示比较完整。它说明了活动背景、经费用途、数量、金额和请求事项，也使用了请示类结尾。学生使用 AI 写请示时，要特别注意金额和审批事项不能编造。若经费需要附明细表，也应在正文中说明附件内容。

\subsubsection*{\textbf{• 函的创作与审查}}

函通常用于联系平级单位或不相隶属单位。学生组织之间需要协调场地、借用设备、邀请协助、报送材料时，都可能使用函的写法。函的语气应正式、平和，既要说明事项，也要提出希望对方配合的内容。

校园 AI 创意作品展示活动结束后，新媒体部准备把活动成果汇编发送给学院团委，并希望学院团委在后续校园科技文化活动中给予指导。小李可以让 AI 起草一份报送函。她先用一个较简单的问题测试 AI。

\begin{userbox}
帮我写一份给学院团委报送活动材料的函。
\end{userbox}

\begin{aibox}
\begin{center}
\textbf{关于报送活动材料的函}
\end{center}

\noindent 学院团委：

\par\qquad 我部现将本次活动相关材料报送给贵单位，请予以审阅。希望贵单位对后续工作给予指导和支持。

\par\qquad 特此函达。

\par\hfill 学生会新媒体部

\par\hfill xx年xx月xx日
\end{aibox}

这份文本虽然具备函的基本样子，但问题也比较明显。它没有写清楚报送的是什么活动材料，也没有说明材料包括哪些内容，更没有体现报送目的。收文单位虽然出现了，但事项过于笼统，阅读者难以判断材料来源和后续处理要求。因此，使用 AI 起草函时，不能只说明写一份函，还应补充收文对象、具体事项、材料名称、希望对方处理的内容、落款和日期等信息。

小李重新补充材料后，再让 AI 起草。

\begin{userbox}
请帮我写一份关于报送校园 AI 创意作品展示活动成果材料的函。收文对象为学院团委。内容包括：活动已结束，现报送活动总结、优秀作品名单和活动照片精选，请学院团委审阅指导。落款为学生会新媒体部，日期为 xx年xx月xx日。
\end{userbox}

\begin{aibox}
\begin{center}
\textbf{关于报送校园 AI 创意作品展示活动成果材料的函}
\end{center}

\noindent 学院团委：

\par\qquad 校园 AI 创意作品展示活动已于近期顺利结束。为便于学院团委了解活动开展情况和成果整理情况，现将活动总结、优秀作品名单和活动照片精选等材料予以报送。

\par\qquad 本次报送材料主要包括：校园 AI 创意作品展示活动总结、优秀作品名单、活动照片精选。请学院团委审阅，并对后续校园科技文化活动开展给予指导。

\par\qquad 特此函达，恳请指导。

\par\hfill 学生会新媒体部

\par\hfill xx年xx月xx日
\end{aibox}

对比前后两次输出可以看到，第二份函更适合正式使用。简单提问得到的文本只保留了函的外壳，内容比较空；补充具体材料后，文本能够明确写出活动名称、报送材料、收文对象和请求事项。通知面向同学，重点是告知任务；请示面向上级，重点是请求批准；函用于联系和报送，重点是说明事项并请求审阅或支持。因此，公文写作不能只看内容，还要判断文种。

\subsubsection*{\textbf{• 会议纪要的整理}}

会议纪要是对会议主要内容和形成事项的整理。活动结束后，小李所在部门召开了复盘会议，讨论了照片整理、优秀作品名单、证书制作、活动总结和下一次活动改进方案。会议纪要不是会议聊天记录的逐字稿，而是要把会议形成的决定和任务分工整理清楚。

小李一开始把会议记录直接交给 AI。

\begin{userbox}
请根据下面内容整理会议纪要。

会议内容：大家讨论了 AI 创意作品展示活动后续工作。部长说照片要赶紧整理。小李说优秀作品名单还要再核对一下。老师提醒总结里要写清楚人数和作品数量。有人说下次现场引导要提前设计路线。还有证书要做，材料要打印。最后说这周内要把总结交上去。
\end{userbox}

\begin{aibox}
\begin{center}
\textbf{AI创意作品展示活动后续工作部署会}
\end{center}

\textbf{会议主题：}AI创意作品展示活动后续工作部署会

\textbf{会议时间：}[待补充]

\textbf{会议地点：}[待补充]

\textbf{主持人：}[待补充]

\textbf{参会人员：}[待补充]

\textbf{一、会议内容}

本次会议围绕AI创意作品展示活动的收尾与总结工作进行了讨论，并形成以下工作部署：

\textbf{活动资料整理}

尽快完成活动照片的筛选与整理工作，确保影像资料及时归档。

\textbf{评选结果复核}

对优秀作品名单进行最终核对，确保获奖信息准确无误。

\textbf{活动总结报告}

严格按照老师的要求，在总结报告中清晰、准确地写明参与人数及作品征集数量。

报告撰写由[待指定]负责，需确保数据真实、内容全面。

\textbf{后续活动优化}

针对下次类似活动，需提前设计并优化现场引导路线，以提升活动组织效率与参与体验。

\textbf{证书制作与材料准备}

即日启动获奖证书的设计与制作工作。

同步完成总结报告及相关材料的打印、装订，做好提交准备。

\textbf{二、会议要求}

以上所有工作须在本周内完成，具体日期待确认，并将最终总结材料按要求上交。

\textbf{三、会议决议}

请相关人员严格按照时间节点推进，确保活动圆满收官。
\end{aibox}

这份纪要已经提炼了主要内容，但缺少会议时间、参会人员和具体分工。小李需要补充会议基本信息和任务安排。

\begin{userbox}
请根据以下信息整理一份会议纪要，要求包含会议时间、会议地点、参会人员、会议内容、形成事项和任务分工。语言正式，重点突出后续任务。

会议名称：校园 AI 创意作品展示活动复盘会议。
会议时间：xx年xx月xx日 19:00。
会议地点：学生会办公室。
参会人员：学生会新媒体部成员及指导老师。
会议内容：讨论活动后续材料整理和下一次活动改进。
形成事项：1. 本周内完成活动总结初稿；2. 两天内完成活动照片筛选和归档；3. xx年xx月xx日前完成优秀作品名单核对；4. 证书制作和材料打印需提交经费请示；5. 下一次活动需提前设计现场参观路线和引导标识。
任务分工：小李负责活动总结初稿和会议纪要；素材整理组负责照片归档；评选协调组负责优秀作品名单核对；部长负责经费请示审核；现场组负责下一次活动引导方案。
\end{userbox}

\begin{aibox}
\begin{center}
\textbf{校园 AI 创意作品展示活动复盘会议纪要}
\end{center}

\textbf{会议时间：}xx年xx月xx日 19:00

\textbf{会议地点：}学生会办公室

\textbf{参会人员：}学生会新媒体部成员及指导老师

\textbf{会议主题：}校园 AI 创意作品展示活动后续材料整理与下一次活动改进

\textbf{一、会议内容}

会议围绕校园 AI 创意作品展示活动后续材料整理、优秀作品名单确认、活动总结撰写、证书制作和下一次活动改进等事项进行了讨论。指导老师提醒，活动总结中应写清楚参与人数、作品数量和活动中存在的问题，避免使用空泛表述。

\textbf{二、形成事项}

1. 本周内完成活动总结初稿。

2. 两天内完成活动照片筛选和归档。

3. xx年xx月xx日前完成优秀作品名单核对。

4. 证书制作和材料打印需提交经费请示。

5. 下一次活动需提前设计现场参观路线和引导标识。

\textbf{三、任务分工}

1. 小李负责活动总结初稿和会议纪要整理。

2. 素材整理组负责活动照片筛选和归档。

3. 评选协调组负责优秀作品名单核对。

4. 部长负责经费请示审核。

5. 现场组负责下一次活动引导方案设计。

\textbf{四、后续要求}

各责任人应按照会议确定的时间节点推进相关工作，完成后及时向部门负责人反馈。涉及名单、经费和对外发布的材料，应在提交或发布前再次核对。
\end{aibox}

第二次纪要更适合实际使用。它不是简单复述会议内容，而是把会议决定转化为任务清单。会议纪要的价值就在于让参会人员知道会后要做什么、由谁负责、什么时候完成。

前面的通知、请示、函和会议纪要示例已经覆盖了常见办公文种，但公文类材料比普通说明文字更需要审查。小李在提交前，要先确认文种、事实、格式、语气和责任是否匹配。

\Needspace{15\baselineskip}
\begin{critiquebox}
公文类材料先审文种。通知、请示、函、纪要用途不同，不能把通知写成请示，也不能把纪要写成总结。公文处理强调实事求是、准确规范、精简高效和安全保密\footnotemark，因此时间、地点、金额、名单、邮箱等事实必须逐项核对。

格式审查要看标题、称呼、正文、结尾、落款和日期是否完整，正式材料还应参照现行公文格式要求检查版式和格式要素\footnotemark。请示还要坚持一事一请，申请经费和报送总结不宜混在同一份请示中\footnotemark。

\textbf{审查要点：}文种是否匹配用途，事实是否能回到原始材料核对，语气是否适合对象，任务是否有责任人和截止时间，未公开名单、联系方式和经费细节是否得到处理。
\end{critiquebox}
\addtocounter{footnote}{-3}
\stepcounter{footnote}\footnotetext{参见：中共中央办公厅、国务院办公厅《党政机关公文处理工作条例》（中办发〔2012〕14号，2012年4月16日发布，2012年7月1日起施行）第五条、第八条。访问日期：2026年7月18日。}
\stepcounter{footnote}\footnotetext{参见：中华人民共和国国家质量监督检验检疫总局、中国国家标准化管理委员会发布的 GB/T 9704—2012《党政机关公文格式》，2012年6月29日发布，2012年7月1日实施。访问日期：2026年7月18日。}
\stepcounter{footnote}\footnotetext{参见：中共中央办公厅、国务院办公厅《党政机关公文处理工作条例》（中办发〔2012〕14号，2012年4月16日发布，2012年7月1日起施行）第十五条。访问日期：2026年7月18日。}



常见公文与会议纪要的用途差异见表\ref{tab:ch4_09_official_document_compare}。

\begin{formalLongTable}
\begin{longtable}{p{0.18\textwidth}p{0.23\textwidth}p{0.23\textwidth}p{0.26\textwidth}}
\caption[常见公文与会议纪要用途比较表]{\chaptertablelabel\quad 常见公文与会议纪要用途比较表}
\label{tab:ch4_09_official_document_compare}\\
\toprule
\textbf{文种} & \textbf{主要用途} & \textbf{常见结尾} & \textbf{审查重点} \\
\midrule
\endfirsthead

\multicolumn{4}{c}{\chaptertablelabel\quad 常见公文与会议纪要用途比较表（续）} \\
\toprule
\textbf{文种} & \textbf{主要用途} & \textbf{常见结尾} & \textbf{审查重点} \\
\midrule
\endhead

\midrule
\multicolumn{4}{r}{续下页} \\
\endfoot

\bottomrule
\endlastfoot

通知 & 告知事项、提出要求 & 特此通知或直接落款 & 对象、任务、时间、反馈方式是否清楚 \\
请示 & 向上级请求批准 & 妥否，请批示 & 是否一事一请，金额和请求事项是否准确 \\
函 & 联系、商洽、报送材料 & 特此函达，盼复或恳请指导 & 收文对象是否合适，事项是否说明清楚 \\
会议纪要 & 记录会议决定和分工 & 一般不写礼貌结尾 & 形成事项、责任人、完成时间是否明确 \\

\end{longtable}
\end{formalLongTable}


通过这些公文和纪要，小李完成了活动后的第三类办公写作任务。可以看到，办公写作不是简单让 AI 写一篇文章，而是围绕真实工作流程进行材料整理、文本生成和人工核对。AI 可以提高整理和起草效率，但不能替代对事实、格式、责任和公开范围的判断。

\subsection*{小结}
本节围绕材料提取、活动总结、方案设计和常用公文展开，强调办公写作要先明确事实、对象和用途，再组织结构和语言。材料提取要保证时间、地点、人数、金额和责任人准确；总结写作要呈现完成情况、实际成效和存在问题；方案设计要写清目标、步骤、分工和保障；公文起草则要注意文种、称谓、格式和行文关系。AI 可以提高整理与起草效率，但最终仍要由使用者核对事实、格式、责任和语气，确保生成内容能够真正用于汇报、通知、协作或存档。
\bigskip

\subsection*{课后习题}
\begin{enumerate}
\setlength{\itemsep}{6pt}

\item 以一次班级活动、社团活动或课程项目为例，列出活动结束后需要整理的原始材料，并按时间、地点、人员、数据、问题和后续任务进行分类。

\item 根据上述材料，设计一段 AI 提问，要求 AI 先提取事实清单，再生成一份简短活动总结，且不能补充原材料中没有的信息。

\item 说明活动总结、活动方案、通知、请示、函和会议纪要在用途、读者对象和内容重点上的区别。

\item 使用 AI 起草办公文档后，应重点核对哪些内容？请从数据金额、名单身份、时间安排、责任分工和发布范围等方面说明。

\item 结合办公写作场景，说明为什么“语言通顺”不等于“文档可用”，并谈谈正式发布或提交前应如何进行人工审查。

\end{enumerate}

\newpage

\section{创意娱乐生成与影像创作}

\begin{sectionkeypoints}
\item 理解创意娱乐生成与影像创作的基本任务，认识小说、漫画和照片修复在文本创作、连续画面和真实影像处理中的不同要求。
\item 掌握 AI 小说生成的基本思路，能够提供人物设定、世界规则、主要冲突和章节目标，避免故事内容空泛或套路化。
\item 了解 AI 漫画分镜生成的处理方法，能够围绕画面顺序、景别、人物动作、背景信息、对白字幕和风格要求组织提示词。
\item 熟悉照片修复的基本边界，能够根据老人旧照片、校园活动照片和证件照等类型区分可修复内容与不可改变内容。
\item 认识创意作品发布和展示的使用边界，明确版权、肖像权、隐私保护、真实性和平台规则等审查要求。
\end{sectionkeypoints}

创意娱乐生成场景，是学生最容易产生兴趣、也最容易看到 AI 直观效果的一类应用场景。它既可以用于写小说、改故事、设计漫画分镜，也可以用于修复旧照片、优化活动照片和整理个人影像资料。与正式材料写作相比，创意娱乐生成的表达空间更大，学生可以加入想象、人物、画面和情绪；但与纯粹娱乐不同，教材中的创意生成仍然要强调方法意识和边界意识。故事可以虚构，真实照片不能被随意伪造；画面可以美化，人物身份和真实场景不能被任意改变；AI 可以帮助生成初稿，但不能替代创作者对内容质量、价值导向和他人权益的判断。

本节围绕小说、漫画和照片修复三个任务展开。小说部分以亲身穿越到古代异世界为连续案例，说明如何从一个模糊想法出发，逐步形成主角身份、世界规则、主要矛盾、章节大纲和正文片段。漫画部分继续使用同一个故事，说明如何把文字剧情拆成连续画面，并处理格数、景别、动作、对白气泡和画面提示词。照片修复部分围绕老人旧照片、校园活动照片和证件照三类常见场景，分别说明可以修复什么、不能改变什么，以及发布前需要检查什么。通过这三个任务，学生可以理解 AI 在创意娱乐中的基本作用：它不是让作品凭空出现，而是帮助使用者整理素材、扩展思路、形成初稿、优化细节，并在人工审查后变成更适合展示和保存的作品。


\subsection{小说灵感生成与文本续写}

小说创作适合用 AI 扩展灵感，但不适合完全交给 AI 自动编写。本小节以亲身穿越到古代异世界为案例，先说明如何把一个模糊想法整理成故事设定，再展示如何让 AI 生成主角、世界观、主要冲突和章节大纲，并说明如何检查人物、情节和语言风格是否前后一致。换言之，AI 在小说创作中的作用，不是替作者决定故事，而是帮助作者把零散想法变成可以继续写下去的文本框架。

\subsubsection*{\textbf{• 案例}}

小李最近想写一篇带有幻想色彩的连载小说。她不想再写普通校园故事，而是希望主角真正进入一个陌生世界：一个看起来像古代王朝，却存在灵力、门派、边境城池和异族商路的世界。主角不是一开始就很强，也没有突然获得无敌能力。他只是一个现代大学生，懂一些现代常识，习惯观察和记录，但在古代异世界里没有身份、没有路引、没有亲人，也不懂当地文字和规矩。这样的故事比单纯开局变强更有张力，因为主角首先面对的不是打败敌人，而是如何在陌生制度中活下来。但如果小李只是简单提问，AI 往往会生成一个比较套路化的开头。

\begin{userbox}
帮我写一个穿越到古代异世界的小说开头。
\end{userbox}

\begin{aibox}
林澈睁开眼时，发现自己躺在一片陌生的荒野上。远处是一座高大的古城，城墙上飘着黑色旗帜，天空中还有几只他从未见过的巨大飞鸟。

他低头看着自己身上的粗布衣服，心中一阵发冷。就在这时，他的手腕忽然亮起一道银色光芒，一个古老的声音在耳边响起：``异乡人，你终于来了。''
\end{aibox}

这个开头虽然有画面感，但问题也很明显。AI 在没有设定依据的情况下，直接安排了黑色旗帜、巨大飞鸟、手腕发光和古老声音。这些元素看起来热闹，却没有解释世界规则，也没有说明主角为什么特殊。如果继续写下去，故事容易变成随机冒险，人物成长也会被神秘力量替代。

\subsubsection*{\textbf{• 内容整理}}

为了让故事更稳定，小李需要先让 AI 整理设定，而不是直接写正文。小李整理出的小说基本设定见表\ref{tab:ch4_11_fantasy_story_setting}。

\begin{formalLongTable}
\begin{longtable}{p{0.24\textwidth}p{0.66\textwidth}}
\caption[古代异世界穿越小说设定整理表]{\chaptertablelabel\quad 古代异世界穿越小说设定整理表}
\label{tab:ch4_11_fantasy_story_setting}\\
\toprule
\textbf{项目} & \textbf{内容} \\
\midrule
\endfirsthead

\multicolumn{2}{c}{\chaptertablelabel\quad 古代异世界穿越小说设定整理表（续）} \\
\toprule
\textbf{项目} & \textbf{内容} \\
\midrule
\endhead

\midrule
\multicolumn{2}{r}{续下页} \\
\endfoot

\bottomrule
\endlastfoot

故事类型 & 古代异世界穿越小说 \\

主角身份 & 现代大学生，意外穿越到陌生古代世界 \\

世界特点 & 有王朝、城池、门派、灵力、边境异族和商路 \\

主角优势 & 会观察、会记录，懂一些现代常识，但没有直接战斗能力 \\

主要困难 & 没有合法身份，不懂文字和规矩，容易被当成流民或奸细 \\

主要冲突 & 主角先解决生存和身份问题，随后卷入边境小城中的失踪案和门派纷争 \\

故事风格 & 有紧张感、探索感和成长感，避免主角开局无敌或缺少合理成长过程的套路化故事 \\

不能改变 & 不要让主角开局无敌，不要突然塞入过多设定，不要让情节全靠巧合推进 \\

\end{longtable}
\end{formalLongTable}

在整理完自己的想法之后，小李继续进行了提问。

\begin{userbox}
我想写一篇古代异世界穿越小说。请不要直接写正文，先帮我整理故事设定。要求如下：

1. 主角是现代大学生，意外穿越到一个类似古代王朝的异世界。
2. 这个世界有城池、门派、灵力和边境异族，但设定不要太复杂。
3. 主角没有开局无敌能力，只能依靠观察、记录、学习和谨慎行动生存。
4. 故事前期发生在一座边境小城，主角先要解决身份问题，再逐渐卷入一件失踪案。
5. 风格要有紧张感和探索感，避免主角开局无敌或缺少合理成长过程的套路化故事。
6. 请输出：故事一句话概括、主角设定、世界观规则、主要矛盾、前三章大纲。
\end{userbox}

\begin{aibox}
\textbf{故事一句话概括}

一名现代大学生意外穿越到有灵力和门派的古代异世界，在边境小城中隐藏身份，学习当地生存的规则，并因一次失踪案逐渐发现这座城背后的秘密。

\textbf{主角设定}

主角名叫林澈，原本是现代大学生。穿越后，他没有强大的武力，也不懂这个世界的文字和规矩。与其他人相比，他的优势是观察细致、习惯记录、善于从细节中推断问题。他的弱点是身体素质普通，对当地制度和危险缺乏了解，因此前期必须谨慎行事。

\textbf{世界观规则}

这个世界类似古代王朝，但存在灵力。普通人生活在城镇和村落中，少数人可以进入门派修习灵术。灵术不能随意使用，城中设有巡检司管理异常事件。边境地区常有异族商队和流民往来，因此身份盘查较严。

\textbf{主要矛盾}

林澈首先要解决身份问题，否则可能被当作流民或奸细。随后，他在客栈帮工时发现一名送药少年失踪。失踪案表面上是普通案件，背后却牵连到城中门派、药铺和巡检司之间的利益冲突。

\textbf{前三章大纲}

第一章：林澈在城外荒坡醒来，发现自己身处陌生世界。他被进城盘查拦下，只能谎称自己是遭遇山匪的外乡人。

第二章：林澈暂时在客栈落脚，通过帮工换取食宿。他开始记录城中的文字、货币、规矩和人物关系。

第三章：一名经常给客栈送药的少年突然失踪。林澈发现少年最后一次出现时，曾提到城北药铺夜里有人搬运木箱。
\end{aibox}

这个回答比直接写正文更适合小说创作。它先确定了故事方向，再确定人物能力和世界规则，避免了主角一开局就获得神秘力量、所有困难都靠巧合解决的问题。之后，小李可以继续要求 AI 生成第一章正文，但需要保留已经确定的设定。

\subsubsection*{\textbf{• 成稿示例}}

在完成故事设定后，小李没有急着扩展后续情节，而是先选择从第一章开头入手。这样做可以检验前面整理出的设定是否能够自然落到正文中，尤其是主角的处境、世界规则和故事氛围是否一致。她进一步向 AI 说明写作范围、字数要求和情节边界，要求 AI 只写主角刚醒来并尝试进城的过程，避免一开始就加入特殊能力或复杂设定。

\begin{userbox}
请根据上面的设定，写第一章开头。要求：

1. 字数 800 字左右。
2. 只写主角醒来、观察环境、尝试进城以及接受盘查的过程。
3. 不要让主角马上获得特殊能力。
4. 重点写陌生感、紧张感和主角的观察。
5. 语言自然，不要堆砌华丽词语。
\end{userbox}

\begin{aibox}
\hspace*{2em}林澈醒来的时候，先听见了风声。

\hspace*{2em}那不是学校宿舍窗外空调外机的声音，也不是图书馆楼下车流经过的声音，而是一种更空、更冷的风声，像从很远的山坡吹过来，贴着地面刮过枯草。

\hspace*{2em}他睁开眼，看到的是灰白色的天空。

\hspace*{2em}林澈愣了几秒，才慢慢撑着地坐起来。身下不是床，也不是地砖，而是一片干硬的土坡。几根枯草扎在袖口边，他低头看了一眼，发现自己身上的卫衣不见了，换成了一件粗布短衣，袖口磨得发白。

\hspace*{2em}他第一反应是自己在做梦。

\hspace*{2em}可是掌心按在地上，石子硌得生疼。风吹进领口，他忍不住打了个寒战。远处有城墙，灰色的墙体横在山坡尽头，城门口排着一队人，有挑担的，有牵驴的，也有背着包袱的。

\hspace*{2em}林澈站在原地，没有立刻往前走。

\hspace*{2em}他先看了看自己的手，又摸了摸口袋。手机没有，学生证没有，钥匙也没有。他身上只有一小块粗布包着的干粮，硬得像石头。

\hspace*{2em}``冷静。''他小声对自己说。

\hspace*{2em}如果这不是梦，那他至少要先弄清楚三件事：这里是什么地方，自己现在是什么身份，怎样才能在天黑前找到可以落脚的地方。

\hspace*{2em}他朝城门走去。越靠近城门，周围的人越多。林澈听不懂所有词，但能听出大概语气。有人在讨价还价，有人在抱怨路难走，还有士兵模样的人站在门口盘查。

\hspace*{2em}轮到林澈时，一个佩刀的守门人上下打量他。

\hspace*{2em}``路引。''

\hspace*{2em}林澈听懂了这两个字，却不知道该怎么回答。他停顿的时间太长，守门人的眼神立刻变了。

\hspace*{2em}``哪来的？''

\hspace*{2em}林澈喉咙发紧。他知道自己不能说实话。一个没有路引、没有口音来源、衣着破旧的人，在边境小城门口解释自己来自另一个世界，只会被当成疯子或者奸细。

\hspace*{2em}他低下头，尽量让声音听起来疲惫一些。

\hspace*{2em}``山路上遇了匪，东西都丢了。我只记得往这边走能进城。''

\hspace*{2em}守门人皱着眉，显然没有完全相信。

\hspace*{2em}林澈站在原地，不敢多说。他发现城门右侧有一张告示，上面贴着几行他看不太懂的字，但旁边画着几张简单的人像，其中一张像是在通缉流民。

\hspace*{2em}他忽然意识到，在这个世界里，不知道规矩本身就是危险。
\end{aibox}

前面的小说片段已经完成了从设定到正文的生成，但连载文本不能只看情节是否热闹。小李需要回到人物设定和世界规则中检查前后是否一致。

\begin{critiquebox}
小说文本审查不只看“好不好看”，更要看设定是否前后一致。林澈如果前面是普通人，后文就不应突然掌握灵术；边境小城、城门盘查和逃亡氛围也要保持同一世界规则。

连载文本还要看情节是否由人物选择和线索推动，而不是依靠巧合强行推动情节或制造转折。语言可以有画面感，但不宜堆砌华丽词语。

\textbf{审查要点：}人物能力是否变化过快，世界规则是否混乱，情节推进是否依赖巧合，语言风格是否稳定，每章结尾是否保留问题或线索。
\end{critiquebox}


\subsection{漫画剧情拆解与分镜设计}

漫画创作不是把小说内容简单配图，而是把故事转化为连续画面。本小节继续使用林澈穿越到古代异世界的案例，先说明如何从小说片段中提取关键画面，再展示如何设计四格或六格漫画分镜，并说明如何检查人物形象、场景氛围、对白气泡和画面风格是否统一。也就是说，漫画生成的重点不是让每一张图都漂亮，而是让读者能按照分格顺序看懂故事。

\subsubsection*{\textbf{• 案例}}

在前一小节中，小李已经使用 AI 完成了小说设定、故事大纲和第一章片段。为了让这个故事更适合展示和传播，她进一步想到，可以把小说中的关键情节改编成漫画。与小说相比，漫画更依赖画面顺序、人物动作和对白安排。小李需要做的，不是让 AI 随意画几张相关图片，而是先把小说内容拆解成连续画面，再让 AI 根据分镜要求生成漫画。

\Needspace{7\baselineskip}
小说第一章中包含多个动作：主角醒来、观察环境、发现城墙、走向城门、被守门人盘查、意识到危险。如果直接让 AI 画成漫画，输出往往会过于笼统，缺少画面顺序和分镜逻辑。更好的做法是进行漫画分镜设计\footnote{漫画分镜设计是将故事情节拆分为连续画面的创作环节，通常需要安排画格顺序、景别变化、人物动作、对白气泡和画面重点，使读者能够按照阅读顺序理解剧情发展。}，先把剧情拆成若干画面，再分别确定景别、人物动作、背景信息和对白。


漫画分镜设计时需要关注的要素见表\ref{tab:ch4_13_comic_storyboard_elements}。

\begin{formalLongTable}
\begin{longtable}{p{0.18\textwidth}p{0.24\textwidth}p{0.46\textwidth}}
\caption[漫画分镜设计要素表]{\chaptertablelabel\quad 漫画分镜设计要素表}
\label{tab:ch4_13_comic_storyboard_elements}\\
\toprule
\textbf{分镜要素} & \textbf{作用} & \textbf{示例} \\
\midrule
\endfirsthead

\multicolumn{3}{c}{\chaptertablelabel\quad 漫画分镜设计要素表（续）} \\
\toprule
\textbf{分镜要素} & \textbf{作用} & \textbf{示例} \\
\midrule
\endhead

\midrule
\multicolumn{3}{r}{续下页} \\
\endfoot

\bottomrule
\endlastfoot

格数 & 决定故事节奏 & 四格适合短情节，六格适合交代环境和转折 \\
景别 & 控制画面信息 & 远景交代城墙，近景表现主角紧张表情 \\
动作 & 推动剧情发展 & 醒来、起身、摸口袋、接受盘查 \\
对白 & 补充人物心理和冲突 & 路引。 我……东西都丢了。 \\
风格 & 保持画面统一 & 古代异世界、写实漫画风、冷色氛围 \\

\end{longtable}
\end{formalLongTable}

如果提问过于简单，AI 只能给出粗略建议。

\begin{userbox}
把林澈穿越到古代异世界的开头改成漫画分镜。
\end{userbox}

\begin{aibox}
第一格：林澈在荒野中醒来，发现自己身处陌生地方。

第二格：林澈看到远处有一座古代城池，感到震惊。

第三格：林澈走到城门口，被守卫盘问身份。

第四格：林澈意识到自己必须隐藏真实来历，故事正式开始。
\end{aibox}

这个回答只能作为初步方向。它没有说明每格画面怎么构图，人物动作是什么，景别如何变化，也没有给出适合图像生成的画面提示词。

\subsubsection*{\textbf{• 内容整理}}

为了让分镜更实用，小李需要提供更明确的要求。

\begin{userbox}
请把下面的小说情节改成六格漫画分镜。要求每一格包含：景别、画面内容、人物动作、对白或字幕、画面提示词。画面采用古代异世界题材的写实漫画风格。主角是一名现代大学生，穿越后身着粗布衣。整体氛围应体现紧张感和陌生感，避免采用夸张、搞笑的画风。

剧情：林澈在城外荒坡醒来，发现自己身处陌生古代世界。他看见远处城墙，摸口袋发现手机和学生证都不见了，只能走向城门。进城时，守门人要求他出示路引。林澈谎称自己在山路上遇到匪，东西都丢了。最后，他意识到不知道这个世界的规矩本身就是危险。
\end{userbox}

\begin{aibox}
\textbf{六格漫画分镜}

第一格，远景。灰白天空下的荒坡，枯草被风吹动，林澈躺在地上刚刚醒来。字幕：他醒来时，听见的不是熟悉的城市声音。提示词：古代异世界荒坡，灰白天空，枯草，少年躺在地上，写实漫画风，冷色氛围。

第二格，中景。林澈撑起身体，低头看着自己身上的粗布衣，表情茫然。对白：这是什么地方？提示词：年轻男生，粗布古代衣服，坐在荒坡上，神情困惑，风吹衣角，细节清晰。

第三格，远景。荒坡尽头出现一座灰色城墙，城门前排着进城的人。林澈站在坡上望向远处。字幕：远处有城墙，也有人群。提示词：古代边境小城，城墙，城门，人群排队，少年背影，空间感强。

第四格，近景。林澈低头摸口袋，发现手机、学生证和钥匙都不见了，手指僵住。旁白：手机没有，学生证没有，连钥匙也没有。提示词：手部近景，粗布衣口袋，少年紧张动作，漫画特写。

第五格，中景。城门口，佩刀守门人拦住林澈，眼神警惕。守门人对白：路引。林澈对白：我……山路上遇了匪，东西都丢了。提示词：古代城门，守门人佩刀，少年低头解释，紧张对峙，写实漫画风。

第六格，特写。林澈看向城门旁的告示，脸色变得凝重。字幕：不知道规矩本身，就是危险。提示词：少年侧脸特写，城门告示，古代文字，紧张氛围，阴影，冷色调。
\end{aibox}

这个分镜比简单提问更适合后续生成漫画画面。它把故事拆成了六个连续画面，有远景、中景、近景和特写，既交代环境，也表现人物心理。对白没有过多解释，而是保留了关键冲突。画面提示词也强调古代异世界、边境小城、粗布衣、冷色氛围等要素，有利于保持画面统一。

\subsubsection*{\textbf{• 成稿示例}}

完成分镜设计后，小李进一步将六格分镜作为画面生成依据，按照已经确定的画格顺序、景别安排、人物动作和场景提示，生成对应的漫画画面。

根据上述分镜生成的漫画示意效果见图\ref{fig:ch4_02_comic}。

\par\vspace{1em}

\par
\noindent\begin{minipage}{\textwidth}
\centering
\BookFigureImage{./images_chapter_04/AI漫画图片.png}
\captionof{figure}{六格漫画分镜示意图}
\label{fig:ch4_02_comic}
\end{minipage}
\par

\par\vspace{1em}

前面的漫画分镜已经把故事拆成连续画面，但分镜能否使用，还要看读者能不能顺着画面读懂故事。小李需要从人物、场景、对白和阅读顺序继续审查。

\begin{critiquebox}
漫画分镜审查要看连续性。林澈的年龄、服装和外貌应保持一致；城外荒坡、远处城墙和城门盘查应属于同一环境，不能一格像中式古城，一格变成欧式城堡。

漫画气泡空间有限，对白要短，小说里的心理描写可以转为字幕或画面动作。六格之间要有动作承接和阅读顺序，不能像六张互不相关的插图。

\textbf{审查要点：}人物形象是否统一，场景风格是否一致，画面顺序是否清楚，对白是否过长，一格是否塞入过多动作或信息。
\end{critiquebox}

使用 AI 生成小说或漫画时，还应注意版权边界。应避免直接照搬他人作品中的人物设定、完整情节、经典对白或具有明显识别性的画面设计；公开发布时，应同时核对版权授权、素材来源和平台对 AI 生成内容的标识要求。

\subsection{照片画面修复与影像优化}

照片修复是创意娱乐生成中最贴近真实生活的一类任务。本小节不把照片修复简单理解为变清晰或变好看，而是根据不同照片类型来确定处理重点：老人旧照片重在保留人物特征和年代感，校园活动照片重在保留现场信息并提高整体清晰度，证件照重在轻度校正而不是改变五官。也就是说，照片修复的第一原则不是美化，而是在尊重真实的前提下改善可读性和保存效果。

很多学生第一次使用 AI 修复照片时，会直接输入把这张照片修好看一点。这种提问看似简单，但风险较大。AI 可能会自动磨皮、年轻化、替换背景、补出不存在的细节，甚至改变人物面部特征。对于风景照片，这可能只是风格变化；但对于老人旧照片、毕业合影或校园活动照片，就可能造成事实失真。因此，学生应先判断照片类型，再说明修复目标和禁止修改的内容。



\subsubsection*{\textbf{• 老人旧照片的修复}}

老人旧照片通常具有纪念意义，修复重点是恢复可辨认度，而不是让照片变成现代写真。比如一张老人年轻时的旧照片，可能存在泛黄、划痕、局部污点和面部模糊。使用 AI 修复时，可以要求去除明显划痕、恢复整体对比度、让人物轮廓更清楚，但不能把老人修成完全不像本人，也不能自动替换服装、重画五官或改变人物原有姿态。

\begin{userbox}
请对这张老人旧照片进行照片修复。重点修复画面中的折痕、划痕、污点和轻微模糊，适当提升人物面部清晰度和整体对比度。请保留黑白旧照片的年代感，不要改成现代写真风格。不要改变人物五官、表情、年龄状态和服装样式。不要过度磨皮，不要让人物年轻化，不要添加原照片中不存在的人物、背景或装饰物。修复后应像一张保存较好的旧照片，而不是重新生成的新照片。
\end{userbox}

老人旧照片修复前后的效果对比见图\ref{fig:ch4_02_old_photo_restore}。

\par\vspace{1em}

\par
\noindent\begin{minipage}{\textwidth}
\centering
\BookFigureImage{./images_chapter_04/老人照片修复.png}
\captionof{figure}{老人旧照片修复前后对比图}
\label{fig:ch4_02_old_photo_restore}
\end{minipage}
\par

\par\vspace{1em}

老人旧照片修复完成后，要特别检查人物特征、年代感和细节真实性。人物应当仍然像原来的人，如果 AI 把脸修得过于光滑、眼睛变大、脸型变窄，就已经不是修复，而是重塑。旧照片可以更清楚，但不应变成高清商业写真。原照片中不存在的细节也不能随意添加，例如原背景只是普通墙面，AI 不应补成华丽室内布景。

\subsubsection*{\textbf{• 校园活动照片的优化}}

校园活动照片常用于公众号、班级总结、社团展示或活动归档，通常涉及多人、横幅、场地、标识和活动事实。修复这类照片时，重点是让画面更清楚、构图更整齐、曝光更正常，而不是改变现场内容。比如讲座合影中横幅较暗、人物略模糊，可以提高亮度和清晰度；但不能把普通班级活动改成校级大型活动，也不能增加不存在的人群或修改横幅文字。


\begin{userbox}
请对这张校园活动照片进行基础修复和优化。重点提升整体亮度、清晰度和对比度，让人物轮廓和后方横幅文字尽量清楚。可以适当裁切画面边缘无关杂物，调整轻微倾斜，使构图更整齐。请不要改变参加人数，不要增加或删除人物，不要修改横幅文字、校名、院名、活动主题和场地信息。不要把普通教室背景替换成礼堂或舞台，不要把活动现场改造成与原照片不一致的场景。修复目标是提高照片可读性，而不是重新生成活动画面。
\end{userbox}



校园活动照片优化前后的效果对比见图\ref{fig:ch4_02_campus_activity_restore}。

\par\vspace{1em}

\par
\noindent\begin{minipage}{\textwidth}
\centering
\BookFigureImage{./images_chapter_04/校园活动照片.png}
\vspace{-0.8em}
\captionof{figure}{校园活动照片优化前后对比图}
\label{fig:ch4_02_campus_activity_restore}
\end{minipage}
\par

\par\vspace{1em}

校园活动照片修复还要注意隐私。如果照片中有人脸、学号、电话号码、定位信息或未经同意展示的个人信息，发布前应进行遮挡或征得同意。AI 优化图片时可能会把模糊文字变得更清楚，这对展示有帮助，但也可能让个人信息暴露得更明显。因此，活动照片修复后的审查不仅要看画面效果，还要看是否适合公开发布。

\subsubsection*{\textbf{• 证件照和个人头像的优化}}

证件照和个人头像的修复边界更严格。它可以调整亮度、背景、衣领和轻微瑕疵，但不应改变五官比例、脸型特征和身份识别信息。比如学生要把一张光线不均匀的头像用于社团报名表，可以让 AI 帮助调整背景和光线；但如果 AI 把脸修得不像本人，或者改变眼睛、鼻子、下巴，就不适合用于正式材料。

\begin{userbox}
请对这张证件照进行轻度修整。可以调整背景颜色，使背景干净统一；可以改善光线不均、轻微阴影和衣领不整齐的问题；可以处理少量临时瑕疵。请不要改变五官比例、脸型轮廓、肤色特征、真实年龄状态和身份识别特征。不要过度磨皮，不要扩大眼睛，不要瘦脸，不要改变发型和表情。修整后应保持自然、正式、真实，适合用于报名表或个人资料。
\end{userbox}


证件照轻度优化的示意效果见图\ref{fig:ch4_02_idphoto_restore}。

\par\vspace{1em}

\par
\noindent\begin{minipage}{\textwidth}
\centering
\BookFigureImage{./images_chapter_04/证件照修复.png}
\captionof{figure}{证件照轻度优化示意图}
\label{fig:ch4_02_idphoto_restore}
\end{minipage}
\par

\par\vspace{1em}

证件照修复的判断标准不是更漂亮，而是更清楚、更正式、更像本人。如果修复后连熟人都觉得不像原来的人，就说明处理过度。对于需要提交到学校、单位或考试系统的照片，还应遵守对应系统的照片要求，不能因为追求美观而造成审核不通过。



前面的照片修复示例展示了不同图片的优化方式，但修复结果不能只看是否更清晰。小李还要判断画面质量提升是否以牺牲真实性为代价。

\begin{critiquebox}
照片修复审查的底线是“改善画面，不改变事实”。老人旧照片要保留人物面貌、年龄状态和原有场景；校园活动照片要核对人数、横幅文字、校名院名和现场环境；证件照或头像可以调整光线，但不能改变五官和身份识别特征。

修复后应对照原图检查人物、场景、文字、数量和关键事实是否发生变化。用于正式材料、活动记录或公开发布时，更要避免过度美化和虚假生成。

\textbf{审查要点：}是否过度年轻化，是否补出不存在的人或物，是否改变活动性质，是否让证件照失真，是否将照片修成与原场景不一致的虚假画面。
\end{critiquebox}


\subsection*{小结}
本节介绍了小说、漫画和照片影像类创作的基本方法。小说创作要先确定人物、情节和世界规则，避免只依赖 AI 随机续写；漫画生成要把故事拆解为连续画面，明确分格、景别、动作和对白；照片修复与影像优化则要以保留真实信息为前提，不能随意改变人物身份和事件背景。AI 可以帮助打开思路、完善分镜和优化画面，但创作过程中仍要保持主题清楚、风格统一，并注意版权、肖像和内容安全。对于公开展示的作品，还应说明生成或修复方式，避免误导读者。
\bigskip

\subsection*{课后习题}
\begin{enumerate}
\setlength{\itemsep}{6pt}

\item 选择一个小说创意，写出主角身份、故事背景、主要冲突和世界规则，并把它整理成一段可直接用于 AI 生成故事大纲的提问。

\item 将一个不少于 200 字的故事片段拆分成 4 至 6 个漫画分镜，写明每格的景别、人物动作、背景信息和对白或字幕。

\item 针对一张老人旧照片、校园活动照片或证件照，说明可以使用 AI 修复哪些内容，以及哪些人物身份、事件背景和真实细节不能被改变。

\item 使用 AI 生成小说、漫画或修复照片后，应如何检查作品的主题是否清楚、风格是否统一、画面是否误导读者？

\item 结合创意娱乐生成场景，说明在公开展示或发布作品前，为什么要审查版权、肖像权、隐私信息和 AI 生成标识。

\end{enumerate}

\newpage

\section{平台内容发布与多平台适配}

\begin{sectionkeypoints}
\item 理解平台内容发布与多平台适配的基本任务，认识真实素材在社交笔记、公众号推文和不同平台版本中的转化要求。
\item 掌握使用 AI 辅助平台发布的基本思路，能够整理地点、人物、事件、照片、感受和注意事项等素材生成标题、正文和图片说明。
\item 了解社交笔记、公众号推文和多平台改写的表达差异，理解个人体验、主题结构、互动表达和语气适配之间的关系。
\item 熟悉 AI 生成内容发布前的审查方法，能够检查内容准确性、标题夸张性、图片授权、肖像信息和表达误导风险。
\item 认识平台发布的使用边界，明确不同平台可以调整形式和语气，但不能改变核心事实、编造体验或误导读者。
\end{sectionkeypoints}

平台发布写作，是把已有素材整理成适合公开传播的内容。学生在日常学习和生活中，经常会遇到类似任务：一次旅行结束后想发社交笔记，一次活动完成后需要写公众号推文，同一件事还可能需要改写成适合朋友圈、班级群等场景发布的内容，或进一步生成短视频标题和公众号摘要。不同平台面对的读者不同，篇幅不同，语气也不同。如果直接把同一段文字复制到所有平台，容易出现表达不合适、重点不清楚、信息过多或过少等问题。

本节以小李的庐山出行为连续案例展开。小李先把一次周末旅行整理成社交笔记，用来记录个人体验并给同学提供参考；随后，她把同一份素材改写成学院公众号推文，用来呈现自然与人文景观观察；之后，她进一步学习如何将同一件事改写成适合不同平台发布的多个版本。通过这一过程，学生可以看到，平台发布不是简单追求热闹文案，而是要在真实素材基础上，根据读者对象和发布目的调整内容。

AI 在平台发布写作中可以发挥辅助作用。它能够帮助学生整理照片信息、提炼标题、生成正文初稿、设计图片说明、调整不同平台语气，也能帮助检查标题是否夸张、结构是否清楚、版本是否一致。但是，AI 不能替学生判断真实经历，也不能替学生承担发布责任。尤其在公开平台发布内容时，学生还需要关注内容准确性、个人信息保护、图片授权、广告与推荐表达、AI 生成内容标识等问题。

因此，本节的学习重点不是记住某个平台的固定模板，而是掌握平台发布的基本方法。学生应从真实素材出发，先明确自己到底想发布什么，再根据平台和读者选择合适结构，避免用同一种表达覆盖所有场景。AI 生成之后，还需要进行人工审查，确认内容真实、语气得体、图片安全、发布合规。只有这样，AI 才能成为平台发布写作的辅助工具，而不是制造模板化、夸张化和失真表达的来源。


\subsection{社交笔记生成与生活分享}

社交笔记是最贴近日常生活的平台内容，常用于记录旅游、展览、探店、校园活动和学习经验。本小节以小李的庐山出行为例，先说明社交笔记可以解决什么表达需求，再介绍如何借助 AI 把零散照片和感受整理成标题、正文和图片说明，并说明发布前应从真实性、图片安全和表达分寸等方面进行审查。换言之，社交笔记不是简单发照片，而是把一次真实经历整理成有信息、有感受、也适合交流的内容。

\subsubsection*{\textbf{• 案例}}

周末，小李和两名同学去了江西著名景点庐山旅游。她拍摄了牯岭街街景、山间步道、云雾景观、观景台以及山路沿线景观的照片，也记下了一些实用信息：从学校出发需要先到九江再转乘景区交通，路程较长，需要提前查询门票和交通信息；山上天气变化较快，热门路段游客较多，傍晚云雾和牯岭街灯光更适合拍照。回来后，她想先发一篇社交笔记，记录这次旅行体验。随后，学院公众号也准备将这次出行经历整理成一篇以自然与人文景观观察为主题的推文。之后，她还需要把同一份素材改成不同平台版本。


小李最初只是想快速得到一篇笔记，于是进行了一个很短的提问。

\begin{userbox}
帮我写一篇庐山旅游笔记。
\end{userbox}

\begin{aibox}
庐山真的很适合周末出游。这里有山间街区、观景台、云雾和夜景，走在山路上会有一种换个环境的感觉。白天可以看风景、拍照片，傍晚可以看云雾和灯光，晚上也很有氛围。

如果你想找一个能看自然风景、也能慢慢走走的地方，庐山是一个不错的选择。建议大家周末约上朋友一起去，提前查好路线和天气，边走边拍照，感受山间风景和街区氛围。
\end{aibox}

这个提问中，小李只提供了庐山旅游笔记这个要求，AI 只能生成比较通用的描述。它不知道小李从哪里出发，不知道她拍了哪些照片，也不知道她真实觉得哪里好、哪里需要避开。因此，这段内容虽然能读，但像许多旅游笔记的通用模板，不适合作为小李自己的笔记直接发布。

\subsubsection*{\textbf{• 内容整理}}

使用 AI 生成社交笔记时，学生要先说明笔记的发布目的和读者对象。社交笔记可以是生活记录，也可以是出行参考，还可以是活动打卡。目的不同，AI 生成内容时选择的重点也不同。如果希望给同学提供出行参考，就不能只写风景好看，还要提供路线、时间、体验感受和注意事项等信息。

同时，要把真实素材整理清楚后再交给 AI。旅游类笔记可以提供出行时间、同行人员、交通方式、游览地点、照片内容、个人感受和需要提醒的地方。对于门票价格、交通信息、排队时间等不确定内容，应要求 AI 不要自行编造，可以提醒读者提前查询。这样生成的笔记既有个人体验，也不会脱离真实经历。

输出要求也要说明清楚。学生可以要求 AI 生成多个标题备选，正文控制在一定字数内，语气自然，不写成广告，并在正文中加入图片配文提示。小李整理好自己的庐山出行经历后，向 AI 提出了更明确的生成要求。


\begin{userbox}
请根据我的真实出行经历，帮我写一篇适合社交平台发布的旅游笔记。要求语气自然，有个人感受，不要写成广告，不要编造没有提供的信息。

我的素材如下：
1. 周末和两名同学去了庐山，从学校出发需要先到九江再转乘景区交通。
2. 主要逛了牯岭街、山间步道、观景台和几处沿线景点，门票和景区交通信息需要提前查询。
3. 我拍了牯岭街街景、山间步道、云雾景观、观景台和山路沿线景观。
4. 热门路段游客比较多，傍晚光线更柔和，云雾和牯岭街灯光更适合拍照。
5. 牯岭街小店人比较多，适合顺路看看，但不建议专门为了某一家店排很久。
6. 我希望这篇笔记能为同学提供周末出行参考，内容控制在 500 字左右。
7. 请给出三个标题备选，并在正文中加入图片配文提示。
\end{userbox}

\begin{aibox}
\textbf{标题备选}

1. 周末出发去庐山，适合和同学慢慢逛

2. 从学校出发的短途小旅行：庐山一日记录

3. 不赶时间的庐山下午，观景台、街灯和山间风

\textbf{正文示例}

这个周末，我和两名同学去了庐山。从学校出发需要先到九江再转乘景区交通，路程需要提前规划，比较适合想周末换个环境、看看自然风景的同学。

门票和景区交通需要提前查询，我们主要逛了牯岭街、观景台、山间步道和几家小店。下午人会比较多，热门路段拍照不太容易避开人群。如果想拍屋檐、观景台和街灯，我更建议傍晚去。那时候光线会柔和一些，牯岭街的颜色也更好看。

这次我比较喜欢的是山间步道。它不像热门路段那么热闹，走起来更安静，也适合和朋友边走边聊天。牯岭街小店的人比较多，我们也排队买了一点，味道还可以，但如果只是为了这家店特意排很久，我觉得没有必要。

图片可以这样安排：

图一：牯岭街街景，适合放在开头，交代庐山氛围。

图二：观景台，适合展示庐山的标志性画面。

图三：山上小吃，可以配一句排队的人很多，味道还可以。

图四：夜景街灯，适合放在结尾，突出傍晚和夜景氛围。

整体来说，庐山不一定适合所有人，但作为周末短途出行还是比较舒服的。它适合愿意提前查好路线、想和同学慢慢走走拍拍的人。
\end{aibox}

第二次输出比第一次更适合发布。它没有编造门票价格、网红店名和所谓必打卡点，而是根据小李提供的真实素材展开。标题也没有使用夸张表达，而是围绕周末短途出行、先到九江再转乘、庐山一日记录等真实信息设计。这样的笔记更可信，也更容易让读者判断是否适合自己。

\subsubsection*{\textbf{• 成稿示例}}

小李在 AI 输出的基础上继续修改。她觉得标题一比较实用，但还可以更自然一点。她也删掉了部分重复表达，把图片配文调整成更像自己会发布的语言。经过修改后，可以形成如下社交笔记成稿。

\begin{aibox}
\textbf{标题：}从学校出发先到九江再转乘，庐山一日小记录

周末和两名同学去了庐山。从学校出发需要先到九江再转乘景区交通，出发前需要查好门票、天气和景区交通信息，适合想在周末换个环境、看看山间风景的时候去走走。

我们主要逛了牯岭街、观景台、山间步道和几家小店。下午人比较多，牯岭街上拍照有点难避开人群。傍晚以后光线会柔和一些，屋檐、观景台和街灯都更适合拍照。

我这次比较喜欢山间步道。热门路段比较热闹，山间会安静一点，适合和朋友边走边聊天。牯岭街小店排队的人不少，味道还可以，但我觉得不用专门为了它排太久。

图片一：牯岭街街景，下午的光线刚好落在瓦片上。

图片二：观景台，路过的人很多，但等一会儿也能拍到比较干净的画面。

图片三：山上小吃，尝一下可以，不用期待太高。

图片四：夜景街灯，傍晚以后比白天更有氛围。

这次出行没有特别紧凑的安排，就是随便走走、拍拍照、聊聊天。如果你也想找一个周末一日出行的地方，庐山可以作为一个备选。
\end{aibox}

这个版本保留了 AI 帮助整理出的结构，也加入了小李自己的判断。比如味道还可以，不用期待太高，随便走走、拍拍照、聊聊天，这些表达更像个人体验，而不是统一模板。社交笔记的特点正在于此：它既要有信息价值，也要保留真实的个人声音。

前面的社交笔记已经把庐山出行整理成可发布内容，但平台发布前不能只看标题是否吸引人。小李还需要检查标题、体验、图片和语气是否会误导读者。

\enlargethispage{4\baselineskip}
\Needspace{15\baselineskip}
\begin{critiquebox}
社交笔记发布前，先看标题和正文是否一致。标题可以吸引人，但不能用正文支撑不了的夸张说法；内容与标题严重不符，容易误导读者\footnotemark。

体验内容也要来自真实经历。没有去过的店、没有核实的交通时间、没有授权的推广内容，都不应写成确定评价或真实测评\footnotemark。图片发布前还要处理同学正脸、车牌、联系方式、票据二维码等信息，个人信息只应在必要范围内使用\footnotemark。

\textbf{审查要点：}标题是否夸张，体验是否真实，图片是否安全，语气是否过度情绪化，评论区回复是否礼貌且不引战。
\end{critiquebox}
\addtocounter{footnote}{-3}
\stepcounter{footnote}\footnotetext{参见：国家互联网信息办公室《网络信息内容生态治理规定》（国家互联网信息办公室令第5号，2019年12月15日公布）第七条第一项。访问日期：2026年7月18日。}
\stepcounter{footnote}\footnotetext{参见：全国人民代表大会常务委员会《中华人民共和国广告法》（2021年4月29日修正）第四条、第二十八条。访问日期：2026年7月18日。}
\stepcounter{footnote}\footnotetext{参见：全国人民代表大会常务委员会《中华人民共和国个人信息保护法》（2021年8月20日公布，2021年11月1日起施行）第五条、第六条。访问日期：2026年7月18日。}



通过这一步，小李完成了一篇适合社交平台发布的旅游笔记。它从个人经历出发，但并不只是随手记录，而是经过了素材整理、AI 初稿、人工修改和发布审查。接下来，同一次旅行还可以被改写成更正式的公众号推文。

\subsection{推文内容策划与图文排版}

公众号推文比社交笔记更强调公共阅读和主题表达，常用于学校、学院、社团和班级发布活动预告、成果展示、经验总结和文化宣传。本小节仍以小李的庐山出行为例，重点说明如何把个人旅行素材转化为适合公众号发布的图文推文：先确定发布主体和读者对象，再设计标题、导语、小标题和图文结构，并对事实边界、图片授权和价值导向进行审查。总体来看，公众号推文不是把社交笔记改得更正式，而是围绕一个明确主题重新组织素材。


\subsubsection*{\textbf{• 案例}}

小李的庐山出行经历原本来自一次周末旅游。后来，学院老师看到小李和同学拍摄的照片，建议她们把这次出行经历整理成一篇公众号推文，主题是从庐山行走中观察自然与人文景观。这个任务和社交笔记不同。它不是单纯告诉同学这里好不好玩，而是要把旅游经历整理成一篇更有主题的图文推送。小李最初没有意识到公众号推文和社交笔记的区别，只是把任务简单交给 AI。

\begin{userbox}
帮我写一篇关于庐山出行的公众号推文。
\end{userbox}

\begin{aibox}
\textbf{标题：}走进庐山，感受自然与人文景观魅力

庐山是一处兼具自然景观和人文积淀的地方。近日，同学们走进庐山，在山间街区中感受自然与人文景观，在山路云雾间体会景区空间的变化。

本次出行中，同学们参观了牯岭街、观景台、山间步道等地点，欣赏了庐山风貌，也记录下了许多精彩瞬间。大家在行走中观察，在交流中思考，对自然与人文景观有了更直观的认识。

通过本次活动，同学们不仅放松了身心，也增强了对本地文化的了解。未来，希望更多同学能够走出校园，在实践中观察社会、认识景区、拓展视野。
\end{aibox}

这段文字比社交笔记正式，但仍然存在明显问题。它把一次普通出行直接写成同学们走进庐山，像是已经组织好的正式活动；它也用了历史韵味、人文气息、景区空间变化等较空泛的表达，却没有说明具体观察内容。由于提问没有提供发布主体、读者对象、图片素材和主题要求，AI 只能生成一篇比较概括的推文。

\subsubsection*{\textbf{• 内容整理}}

使用 AI 写公众号推文时，学生需要先把发布场景说清楚。公众号推文比普通社交笔记更强调主题和结构，因此要让 AI 知道文章发布在哪类公众号上，面向哪些读者，想表达什么主题，语气是自然正式还是轻松活泼。

同时，要把真实素材整理好。可以提供事件背景、主要经过、观察内容、照片素材和想表达的收获，也要说明哪些内容不能写。比如个人经历不能被写成官方活动，没有发生的讲话、人数、成果和评价不能让 AI 自动补充。生成要求也要写清楚，例如要求 AI 先给出标题、导语、正文小标题、各部分写作要点和图片安排建议。这样得到的内容更像一份可修改的推文策划稿，而不是直接生成一篇难以调整的完整文章。

\begin{userbox}
请根据以下真实素材，帮我策划一篇适合学院公众号发布的推文。主题是从庐山一日行中观察自然与人文景观。请不要把它写成官方组织的大型活动，也不要编造没有发生的内容。

真实素材如下：
1. 小李和两名同学周末自行前往庐山，从学校出发需要先到九江再转乘景区交通。
2. 主要观察了牯岭街街景、山间步道、云雾景观、观景台和山路沿线景观。
3. 出行目的原本是旅游记录，后来希望整理成一次自然与人文景观观察分享。
4. 可以表达的收获包括：庐山适合观察自然景观、山地街区、游客活动和景区公共服务；傍晚云雾和牯岭街灯光让景区氛围更明显；山间步道体现了休闲空间的作用。
5. 推文面向学院公众号读者，语气要自然正式，适合学生阅读。
6. 请给出标题、导语、正文小标题、每部分写作要点和图片安排建议。
\end{userbox}

\begin{aibox}
\textbf{标题建议}

1. 从庐山出发，记录一次自然与人文景观观察

2. 庐山一日行：在牯岭街、观景台与灯火中观察景区

3. 走出校园，在庐山感受景区空间

\textbf{导语示例}

一个周末的下午，小李和两名同学从学校出发来到庐山。最初，这只是一次普通的短途出行；整理照片和记录时，她们发现，牯岭街、观景台、山间步道和夜景街灯中，也藏着可以观察自然与人文景观的细节。

\textbf{正文结构建议}

\textbf{一、从校园到庐山：一次自然与人文景观观察}

写作要点：交代目的地、路程时间和出行背景，说明这不是大型实践活动，而是一次从日常出行延伸出来的观察。

图片建议：公交站或庐山入口照片，配文为从学校出发需先到九江再转乘景区交通，短途出行也可以成为观察景区的开始。

\textbf{二、牯岭街与观景台：山地景区空间中的生活痕迹}

写作要点：结合牯岭街街景、观景台和沿线小店，写出山地街区如何承载游客行走、拍照和消费活动。

图片建议：牯岭街街景、观景台照片，配文为牯岭街和观景台让景区的空间层次更明显。

\textbf{三、山间步道与夜景街灯：休闲空间的另一面}

写作要点：写山间步道比热门路段安静，适合停留和交流；傍晚灯光改变了景区氛围，也让游客体验发生变化。

图片建议：山间步道和夜景街灯照片，配文为傍晚以后的庐山更适合慢慢观察。

\textbf{四、从旅游记录到文化观察}

写作要点：收束全文，强调日常出行也可以成为观察景区、理解文化和训练表达能力的机会。

图片建议：三人背影或远景照片，避免出现清晰正脸。
\end{aibox}

第二次输出不是直接给出一篇完整推文，而是先形成了策划方案。公众号写作中，这一步很重要。许多学生写推文时直接进入正文，结果容易造成标题、导语、正文和图片彼此脱节。先用 AI 生成结构和图文安排，再进行人工判断，可以减少后续修改成本。

% \subsubsection*{\textbf{• 成稿示例}}
% 公众号推文成稿示例已合并至正文。
在确认结构后，小李可以继续让 AI 根据方案生成完整推文。此时，AI 已经有了标题、导语、分层结构和图片安排，输出就会更接近正式成稿。

\begin{userbox}
请根据刚才的推文策划，生成一篇完整的学院公众号推文初稿。要求保留真实边界：这是小李和两名同学的周末自行出行，不是学院统一组织的大型活动。推文主题是从庐山一日行中观察自然与人文景观。语言自然正式，不要夸大意义，不要编造采访、老师讲话或集体活动环节。正文控制在 1200 字左右。
\end{userbox}

\begin{aibox}
\textbf{从庐山出发，记录一次自然与人文景观观察}

\textbf{导语}

一个周末的下午，小李和两名同学从学校出发来到庐山。最初，这只是一次普通的短途出行。整理照片和记录时，她们发现，牯岭街、观景台、山间步道和夜景街灯中，也藏着可以观察自然与人文景观的细节。

\textbf{一、从校园到庐山：一次自然与人文景观观察}

从学校到庐山，需要先到九江再转乘景区交通。这样的行程需要提前规划，也让一次周末出行更适合做成有准备的观察记录。对大学生来说，自然与人文景观观察并不一定只发生在正式课程或大型实践中。有时，一次普通的出行、一次停留、一次拍照和一次记录，也可以成为理解景区空间的开始。

这次出行中，小李和同学没有安排复杂路线，而是沿着牯岭街慢慢行走。她们把看到的屋檐、观景台、店铺、街灯和山间步道记录下来，也在整理照片时思考这些空间为什么会吸引游客停留。

图片建议：庐山入口或公交到达后的第一张照片。图片说明：从学校出发需先到九江再转乘景区交通，一次短途出行也可以成为观察景区的开始。

\textbf{二、牯岭街与观景台：山地景区空间中的生活痕迹}

庐山的牯岭街和观景台，是这次出行中最容易被注意到的画面。屋檐、石板路和临街小店构成了山地景区的基本空间，也让游客在行走中形成对庐山的第一印象。

在牯岭街上，游客拍照、沿路停留和购买小吃。山地景区空间在这里并不是静止的背景，而是和当下的消费、休闲、社交活动结合在一起。小李发现，很多人并不会完整了解庐山相关历史，但会通过拍照、品尝小吃和沿街行走形成自己的体验。

图片建议：牯岭街街景和观景台照片。图片说明：牯岭街和观景台让景区的空间层次更明显，也让游客的行走路线更有停留感。

\textbf{三、山间步道与夜景街灯：休闲空间的另一面}

相比热门路段，山间步道更安静一些。这里没有太多店铺，更多是散步、聊天和停留的人。对小李来说，山间步道是这次出行中比较舒服的一段，因为它不像热门路段那样拥挤，也更适合观察游客如何使用公共休闲空间。

傍晚以后，街灯亮起，景区的氛围发生了变化。白天更适合观察建筑和街道，傍晚则更适合感受光线和人流带来的变化。小李在整理照片时发现，同一个地点在不同时间会呈现出不同效果，这也是景区空间观察中值得注意的细节。

图片建议：山间步道和夜景街灯照片。图片说明：傍晚以后的庐山更适合慢慢观察，也更容易感受到空间氛围的变化。

\textbf{四、从旅游记录到文化观察}

这次庐山一日行并不是一次正式调研，但它提供了一个提醒：自然与人文景观并不只存在于书本和展馆中，也存在于山路、街区、观景台、云雾和人的活动里。对大学生来说，日常出行也可以成为观察景区、训练表达和积累素材的机会。

从一篇社交笔记到一篇公众号推文，内容的重点发生了变化。前者更重视个人体验，后者更重视主题表达和公共阅读。借助 AI，小李可以更快地完成结构整理和语言初稿；但素材是否真实、主题是否合适、图片是否对应，仍需要自己认真判断。

图片建议：庐山远景或背影照片。图片说明：走出校园，日常行走也可以成为理解自然与人文景观的一种方式。
\end{aibox}

这个公众号初稿已经具备标题、导语、分层正文和图片说明。它没有把小李的个人旅游夸大成学院组织的大型实践，也没有编造老师讲话、采访内容和活动成果。因此，它比第一次简短提问得到的推文更符合真实材料。

\subsubsection*{\textbf{• 公众号推文的图文排版}}

公众号推文不仅是文字，还包括阅读节奏。图片放在哪里、图片说明怎么写、小标题是否清楚，都会影响阅读效果。对学生组织和学院公众号来说，排版不一定要复杂，但要做到图文对应、层次清楚、信息准确。


公众号推文的图文排版安排见表\ref{tab:ch4_18_wechat_layout}。

\begin{formalLongTable}
\begin{longtable}{p{0.2\textwidth}p{0.36\textwidth}p{0.34\textwidth}}
\caption[公众号推文图文排版安排表]{\chaptertablelabel\quad 公众号推文图文排版安排表}
\label{tab:ch4_18_wechat_layout}\\
\toprule
\textbf{位置} & \textbf{图片内容} & \textbf{作用} \\
\midrule
\endfirsthead

\multicolumn{3}{c}{\chaptertablelabel\quad 公众号推文图文排版安排表（续）} \\
\toprule
\textbf{位置} & \textbf{图片内容} & \textbf{作用} \\
\midrule
\endhead

\midrule
\multicolumn{3}{r}{续下页} \\
\endfoot

\bottomrule
\endlastfoot

封面图 & 庐山入口、观景台或夜景街灯 & 建立整体氛围，吸引读者进入正文 \\

导语后 & 庐山入口或出发路程照片 & 交代出行背景和地点 \\

第一部分 & 牯岭街街景、观景台 & 展示山地街区和空间层次 \\

第二部分 & 牯岭街小店、沿线小店 & 展示游客活动和商业空间 \\

第三部分 & 山间步道、夜景街灯 & 展示休闲空间和时间变化 \\

结尾处 & 远景、背影或合照剪影 & 收束全文，避免过多展示清晰正脸 \\

\end{longtable}
\end{formalLongTable}

使用 AI 进行排版建议时，可以让 AI 输出图片位置、图片说明和对应段落。这样做的好处是，学生不必在写完正文后才临时找图，而是从策划阶段就把图文关系考虑进去。公众号图文写作应避免图片与文字相互脱节，例如文字讲山间步道，图片却放小吃；文字讲文化观察，图片却全部是自拍。AI 可以帮助提醒这种对应关系，但最终仍需要发布者自己核对照片内容。

前面的公众号推文已经完成主题化组织和图文排版，但公众号内容的传播范围更大。小李在发布前，需要把个人体验、学院账号边界、图片授权和事实信息重新核对。

\Needspace{12\baselineskip}
\begin{critiquebox}
公众号推文要先审发布主体。个人账号可以保留个人感受，学院或学生组织账号则不能把个人体验写成正式立场，也不能把自行出行写成学院组织活动。公众账号内容应保持积极、健康、负责的表达边界\footnotemark。

事实信息要逐项核对，尤其是时间、地点、人数、路线、活动性质、参与人员和图片说明。AI 可能自动补充没有依据的内容，生成后的文本必须删除无来源信息，保留准确可靠的表述\footnotemark。

\textbf{审查要点：}发布主体是否清楚，事实信息是否核实，图片人像是否授权，标题导语是否稳妥，表达是否自然正式、积极健康。
\end{critiquebox}
\addtocounter{footnote}{-2}
\stepcounter{footnote}\footnotetext{参见：国家互联网信息办公室《互联网用户公众账号信息服务管理规定》（2021年1月22日公布）第四条。访问日期：2026年7月18日。}
\stepcounter{footnote}\footnotetext{参见：国家互联网信息办公室、国家发展和改革委员会、教育部等《生成式人工智能服务管理暂行办法》（2023年7月10日公布）第四条第五项。访问日期：2026年7月18日。}



公众号推文完成后，小李发现，同一份素材在不同平台上的表达差异很大。社交笔记重在个人体验，公众号推文重在主题表达。接下来，她还需要把同一件事改写成适合不同平台发布的版本。


\subsection{多平台内容改写与表达适配}

同一件事情可以发布到不同平台，但不同平台面对的读者、篇幅、语气和表达目的并不相同。本小节继续以小李的庐山出行为例，说明如何把同一份素材改写成社交笔记、公众号摘要、朋友圈内容和班级群分享：先建立素材卡，明确不能改变的事实边界；再让 AI 生成不同平台版本；并检查各版本的事实是否一致、语气是否合适、标题是否过度夸张。总的来说，多平台改写不是简单换几个词，而是在不改变事实的前提下重新安排表达重点。


AI 很适合辅助多平台改写。它可以把一份原文改成不同长度、不同语气、不同结构，也可以帮助生成标题、摘要、图片说明和互动回复。但学生必须告诉 AI：哪些信息必须保留，哪些信息可以调整，哪些内容不能编造。否则，AI 可能为了迎合平台风格加入夸张表达，甚至把普通旅游写成爆火攻略或官方活动。

\subsubsection*{\textbf{• 案例}}

以小李的庐山出行为例，如果发社交笔记，重点是路线、体验和照片；如果发学院公众号，重点是自然与人文景观观察；如果发朋友圈，重点是个人感受；如果发班级群，重点是给同学提供出行参考；如果发短视频标题，重点是用简短句子概括画面亮点。同一份素材可以产生多个版本，但每个版本都要保留核心事实，不能为了平台风格而改变事情本身。

小李先根据经历尝试了一个简单提问。

\begin{userbox}
周末我和两位朋友去了庐山。主要逛了牯岭街、山间步道、观景台和几处沿线景点，门票和景区交通信息需要提前查询。热门路段游客比较多，傍晚云雾和光线更适合拍照。牯岭街小店人比较多，适合顺路看看，但不建议专门为了某一家店排很久。我们拍了牯岭街街景、山间步道、云雾景观、观景台和山路沿线景观。请根据以上素材生成社交平台版、公众号版、朋友圈版、班级群分享版文案。
\end{userbox}

\begin{aibox}
\textbf{1. 社交平台版}

标题：庐山周末漫游｜不赶景点，只拾云雾与山风

和两位朋友躲进庐山过周末，没有按头打卡，主打一个随性溜达。

路线：牯岭街闲逛 → 山间步道穿行 → 观景台等雾散 → 沿线随停随拍

Tips：

门票和景区交通记得提前线上查好，现场排队费时。

热门路段人从众，建议早出发或午后错峰。

傍晚时分绝了！光线变柔，云雾流动，观景台和山路都像加了滤镜。

牯岭街小店氛围不错，但人挤人，顺路看看就好，不必为某家店排长队。

出片点：街景老房子、步道树影、云雾翻涌、山路弯道——我们都拍了，九宫格放不下。

总结：庐山适合慢慢走，别赶路，去感受光与雾的即兴演出。

\textbf{2. 公众号版}

\begin{center}
\textbf{三人庐山周末：在云雾里，找回慢走的节奏}
\end{center}

\par\qquad 周末，与两位朋友相约庐山。没有制定严密的行程表，只定了一个方向——上山，走走。

\par\qquad 我们从牯岭街开始，这条云端上的街市比想象中热闹。沿街小店琳琅满目，但游客络绎不绝，我们默契地达成共识：顺路看看即可，不值得为任何一家店耗费太久排队。真正的风景，在街巷之外。

\par\qquad 转入山间步道，人声渐远。沿途经过几处观景台，恰好遇上傍晚时分——这是庐山赠予我们最好的时刻。光线从硬朗变得柔软，云雾从谷底升腾又散开，山路的曲线在暮色中格外动人。我们拍了街景的怀旧感、步道的幽静、云雾的变幻、观景台的辽阔，以及山路蜿蜒的线条，每一帧都像是天然滤镜。

\par\qquad 给想去的朋友一点提醒：门票和景区交通车信息变动快，出发前一定在官方渠道确认最新价格和运营时间。热门观景点和步道入口上午人流量大，若想安静拍照，可考虑午后至傍晚时段。牯岭街的餐饮和小店适合路过感受，不必抱有``必吃必买''的执念。

\par\qquad 庐山的美，不在打卡清单，而在你停下脚步时，恰好有一阵雾路过。

\textbf{3. 朋友圈版}

\par\qquad 周末庐山，三人行。走了牯岭街，钻了山间步道，蹲了几个观景台，等了一场傍晚的云雾和光。人确实不少，但好在风景没掉线——傍晚那会儿，光线软得像调过色，山路转弯处随手一拍都好看。小店逛了，人多，没凑热闹排队，顺路看看就挺好。

\par\qquad 门票交通记得提前查，别的随缘。相册里多了几十张云雾和山影，这个周末，值了。

\textbf{4. 班级群分享版}

@所有人 周末庐山小分队汇报！

我和两个朋友去庐山晃了两天，路线很佛系：牯岭街→步道→观景台→沿线乱走。

重点说几个给以后想去的同学参考：

门票和景区车信息一定要提前查，我们差点扑空。

热门路段人巨多，建议傍晚出动，光线绝美，云雾也配合，拍照直接原图出。

牯岭街小店逛逛可以，但别为某家店排长队，真没必要。

我们拍了街景、步道、云雾、观景台和山路弯道，照片等我整理好发群里！谁下次想去可以喊我，我还能再冲一次，但一定要选傍晚拍照哈哈哈哈。
\end{aibox}

这个回答看起来给出了多个平台版本，但仍然存在两个问题：不同平台之间的内容边界感不强，每个部分都有攻略部分；同时，由于原始经历的口语化叙述较多，生成文本也保留了明显的口语化问题。

\subsubsection*{\textbf{• 内容整理}}

多平台改写前，可以先建立一张素材卡。小李根据之前学到的材料整理的知识，使用AI先按照地点、过程、照片、体验和注意事项等项目，把庐山出行整理成了以下素材卡。


小李整理出的庐山出行素材卡见表\ref{tab:ch4_20_lushan_material_card}。

\begin{formalLongTable}
\begin{longtable}{p{0.24\textwidth}p{0.66\textwidth}}
\caption[庐山出行多平台改写素材卡表]{\chaptertablelabel\quad 庐山出行多平台改写素材卡表}
\label{tab:ch4_20_lushan_material_card}\\
\toprule
\textbf{项目} & \textbf{内容} \\
\midrule
\endfirsthead

\multicolumn{2}{c}{\chaptertablelabel\quad 庐山出行多平台改写素材卡表（续）} \\
\toprule
\textbf{项目} & \textbf{内容} \\
\midrule
\endhead

\midrule
\multicolumn{2}{r}{续下页} \\
\endfoot

\bottomrule
\endlastfoot

事件 & 三人周末前往庐山 \\
地点 & 庐山、牯岭街、山间步道、观景台 \\
费用 & 门票和景区交通需提前查询，其他消费按个人实际情况 \\
素材 & 牯岭街街景、山间步道、云雾景观、观景台和山路沿线景观照片 \\
真实体验 & 热门路段游客较多，傍晚云雾和光线更适合拍照，牯岭街小店人较多，山间步道体验慢节奏生活 \\
不能改变 & 不能写成学院统一组织活动，不能编造店名和价格，不能夸大为必去景点 \\

\end{longtable}
\end{formalLongTable}

小李随后使用素材卡重新提问。

\begin{userbox}
请根据素材卡，把同一件事改写成四个不同平台版本。要求每个版本都保留核心事实，不要编造，不要把个人出行写成学院组织活动。

请分别生成：
1. 社交笔记版，300 字左右，偏旅游攻略。
2. 学院公众号版，300 字左右，主题是自然与人文景观观察，语气自然正式。
3. 朋友圈版，80 字左右，偏个人感受与生活分享。
4. 班级群分享版，100 字左右，偏实用提醒。
\end{userbox}

\begin{aibox}
\textbf{一、社交笔记版}

庐山周末漫游｜三人随性走，傍晚光影最出片

周末和两位朋友上山，没有赶景点，只走了几条经典路线：牯岭街、山间步道、沿途观景台。真实体验分享给想去的你——

行前必看：门票和景区交通车信息变动快，出发前一定在官方渠道确认最新价格和运营时间，现场再查容易耽误。

路线感受：

牯岭街热闹，小店多但人也多，顺路逛逛氛围不错，不必为某家店排长队。

转入山间步道，速度瞬间慢下来，适合体验慢节奏生活。

观景台和山路沿线视野开阔，但热门路段游客扎堆，建议下午三四点后再去。

拍照黄金时段：傍晚云雾流动，光线变柔，山脊和街道都像镀了层暖色。我们拍的街景、步道、云雾、观景台、山路弯道，全是在那个时段最出片。

花费：除门票交通外，餐饮购物丰俭由人，我们人均没超预期。

总结：庐山宜慢逛，别信必去标签，随遇而安反而惊喜更多。

\textbf{二、学院公众号版}

\begin{center}
\textbf{庐山周末一瞥：山间步道与街市光影}
\end{center}

\par\qquad 周末与两位友人同往庐山，以个人游历者的视角，记录下这处自然与人文交织的山间图景。

\par\qquad 牯岭街是庐山生活气息的集中体现。街面店铺鳞次栉比，游人往来不绝，热闹中透出小镇特有的烟火气。然而过于密集的人流也让深入探店变得不易。我们仅是顺路观望，未作长时间停留，反倒将更多时间留给山野。

\par\qquad 转至山间步道，景致骤然切换。石阶蜿蜒，林木掩映，偶有观景台伸出崖壁，可俯瞰层峦叠嶂。此时游客渐少，脚步也随之放缓，真正体会到慢节奏的意味。傍晚时分的云雾与斜阳构成最动人的画面：光线不再刺眼，雾气时而聚拢时而飘散，山路沿线处处可入镜。我们拍摄了街景、步道、云雾、观景台与山路，这些影像如实记录了光线与地形共同塑造的瞬间。

\par\qquad 此行也观察到，热门观景台在白天人流较多。若想安静感受自然，不妨选择午后至傍晚时段。门票与交通信息需提前查询，避免现场被动。整体而言，庐山的美在于其多变的光影与松弛的步行体验，而非某个被标注的必到之处。

\textbf{三、朋友圈版}

\par\qquad 周末和俩朋友上庐山，不赶景点，只走了牯岭街、步道和观景台。傍晚那会儿云雾和光线绝了，随手拍都好看。小店人多，没去凑热闹，步道慢慢走很舒服。门票交通记得提前查，相册已满，心还留在山上。

\textbf{四、班级群分享版}

\par\qquad @所有人 周末跟两个朋友去了庐山，说几个实用点：门票和景区车一定要提前网上查好，现场容易懵；牯岭街小店人多，别为某家排长队，不划算；下午到傍晚去观景台和步道人少光线好，拍照最佳；步道适合慢走聊天，不用赶。我们拍了些照片，回头发群里，想去的可以参考。
\end{aibox}

第二次输出更能体现不同平台之间的差异化表达。不同版本的语气和重点已经有所区别，但核心事实没有改变。社交笔记更像旅游攻略，公众号版强调自然与人文景观观察，朋友圈版注重体验，班级群分享版强调实用提醒，边界感更强。同时，部分正式文案口语化问题得到缓解，例如公众号版。

\subsubsection*{\textbf{• 平台版本差异比较}}

为了帮助学生理解不同平台的表达差异，可以把同一份素材拆成几个维度来比较。平台不是简单的发布位置，而是读者、目的、语气和篇幅的组合。


多平台内容版本之间的差异见表\ref{tab:ch4_21_platform_version_compare}。

\begin{formalLongTable}
\begin{longtable}{p{0.16\textwidth}p{0.22\textwidth}p{0.22\textwidth}p{0.28\textwidth}}
\caption[多平台内容版本差异比较表]{\chaptertablelabel\quad 多平台内容版本差异比较表}
\label{tab:ch4_21_platform_version_compare}\\
\toprule
\textbf{平台版本} & \textbf{主要读者} & \textbf{表达重点} & \textbf{写作特点} \\
\midrule
\endfirsthead

\multicolumn{4}{c}{\chaptertablelabel\quad 多平台内容版本差异比较表（续）} \\
\toprule
\textbf{平台版本} & \textbf{主要读者} & \textbf{表达重点} & \textbf{写作特点} \\
\midrule
\endhead

\midrule
\multicolumn{4}{r}{续下页} \\
\endfoot

\bottomrule
\endlastfoot

社交笔记 & 对旅游和生活分享感兴趣的读者 & 个人体验、路线、照片和实用提醒 & 标题自然，正文有体验感，可加入图片配文 \\

公众号推文 & 学院师生、关注活动或主题的读者 & 主题表达、观察收获、图文结构 & 导语清楚，小标题分层，语气自然正式 \\

朋友圈 & 熟人关系中的同学和朋友 & 心情记录和生活片段 & 篇幅短，语气轻松，不需要完整攻略 \\

班级群分享 & 可能想出行的同学 & 时间、交通、费用和注意事项 & 信息直接，少写修饰，方便同学判断 \\

\end{longtable}
\end{formalLongTable}

这个表格说明，同一件事情在不同平台上不是换几个词就可以。比如社交笔记可以写傍晚更适合拍照，但公众号推文应把它转化为空间氛围变化；朋友圈可以写走走拍拍很舒服，但班级群更适合写交通、门票和景区交通需提前查询、下午人多等信息。

前面的多平台改写已经生成不同版本，但版本越多，越容易出现事实不一致和语气失控。小李需要先建立同一份素材基准，再判断各平台表达是否适配。

\Needspace{12\baselineskip}
\begin{critiquebox}
多平台改写首先要保证事实一致。社交笔记、公众号、朋友圈和班级群可以语气不同，但所有版本都应以同一份素材卡为准。

其次要控制平台化夸张。轻松表达可以有，但不应使用“后悔才来”“绝美到失语”等与事实不匹配的强情绪标题；涉及店铺、商品或服务评价时，也不能把未经核实的体验写成推荐或宣传\footnotemark。若大量使用 AI 生成文本、图片、视频或声音，还应按平台规则和相关要求进行说明或标识\footnotemark。

\textbf{审查要点：}核心事实是否一致，语气是否适配平台，标题是否过度夸张，商业表达是否有边界，隐私信息是否扩散，AI 生成内容是否需要标识。
\end{critiquebox}
\addtocounter{footnote}{-2}
\stepcounter{footnote}\footnotetext{参见：全国人民代表大会常务委员会《中华人民共和国广告法》（2021年4月29日修正）第四条、第二十八条。访问日期：2026年7月18日。}
\stepcounter{footnote}\footnotetext{参见：国家互联网信息办公室、工业和信息化部、公安部、国家广播电视总局《人工智能生成合成内容标识办法》（国信办通字〔2025〕2号，2025年3月7日发布）第十条、第十四条。访问日期：2026年7月18日。}



\subsection*{小结}
本节围绕社交笔记、公众号推文和多平台改写展开，说明同一内容在不同平台上需要调整标题、结构、语气和呈现方式。社交笔记更强调体验感和互动性，公众号推文更强调主题策划、段落层次和图文配合，多平台适配则要根据平台长度、读者习惯和传播目的进行重写。AI 的作用是辅助改写与适配，而不是替代发布者对事实和影响的判断。发布前仍要检查真实性、边界意识和平台规范，尤其要避免夸张标题、虚假体验、未经授权的图片使用，以及涉及他人隐私的信息泄露。


\bigskip
\subsection*{课后习题}
\begin{enumerate}
\setlength{\itemsep}{6pt}

\item 选择一次真实旅行、展览、校园活动或学习经历，整理可用于平台发布的素材，包括地点、人物、事件、照片内容、个人感受和注意事项。

\item 根据同一份素材，分别设计一段用于生成社交笔记和公众号推文的 AI 提问，要求体现读者对象、篇幅要求、语气差异和图片说明。

\item 对比社交笔记、公众号推文、班级群分享和朋友圈短文，说明它们在表达重点、语言风格和信息密度上的不同。

\item 使用 AI 改写平台内容后，应如何检查标题是否夸张、事实是否准确、图片是否获得授权、他人隐私是否被泄露？

\item 结合平台内容发布场景，说明为什么多平台适配可以改变表达形式，却不能改变核心事实，并谈谈发布者应承担哪些审查责任。

\end{enumerate}

\section*{全章回顾}
\addcontentsline{toc}{section}{全章回顾}

本章围绕人工智能在通识应用场景中的使用方法展开，重点说明学生如何在正式材料写作、办公文档生成、创意娱乐创作和平台内容发布中合理使用 AI。无论是生成简历、申请信、活动总结，还是创作小说、漫画、修复照片和改写平台内容，都应先明确任务目标，再整理真实素材，生成后对 AI 输出进行人工审查。

通过本章学习可以看到，AI 的作用不是替代人的判断，而是帮助学生完成材料整理、结构生成、表达优化和版本适配。使用 AI 时，学生既要提高提问的清晰度，也要关注事实真实性、隐私保护、版权意识和内容标识等问题。只有把 AI 生成结果作为可修改的初稿，而不是直接发布或提交的成品，才能更安全、有效地发挥人工智能工具的辅助价值。
